\chapter{Riemannian Metrics}
\label{chap:pet-riemannian-metrics}

Faithful transcription of Petersen, \emph{Riemannian Geometry} (GTM 171, 3rd ed.)~\cite{petersen2016},
Chapter 1. The chapter introduces Riemannian manifolds and the basic maps between
them (isometries, isometric immersions, Riemannian submersions), the volume form,
the role of isometry groups and Lie groups in constructing Riemannian manifolds,
local (coordinate and frame) representations of a metric — including the family of
rotationally symmetric and doubly warped product metrics that produce the space
forms and the Hopf fibrations — and the tensor-algebra conventions (type change,
contraction, the Euclidean inner product of tensors) used throughout the book. Each
numbered statement keeps Petersen's example/proposition number in the environment
title; proofs keep his explicit cross-references so the dependency graph is
recoverable from the text.

\section{Riemannian Manifolds and Maps}

\begin{definition}[Riemannian manifold]
  \label{def:pet-ch1-riemannian-manifold}
  \lean{PetersenLib.RiemannianMetric}
  \leanok
  A \emph{Riemannian manifold} $(M,g)$ consists of a $C^\infty$-manifold $M$
  (Hausdorff and second countable) together with a Euclidean inner product $g_p$
  (also written $g|_p$) on each tangent space $T_pM$, such that $p\mapsto g_p$
  varies smoothly: for any two smooth vector fields $X,Y$ the function
  $p\mapsto g_p(X|_p,Y|_p)$ is smooth. The subscript $p$ is usually suppressed, so
  one writes $g(X,Y)$, understood pointwise; the metric is also denoted $g_M$ when
  it needs to be tied to $M$. The tensor $g$ is the \emph{Riemannian metric}, or
  simply the \emph{metric}. The manifold is generally assumed connected, with
  occasional exceptions (level sets, submanifolds cut out by constraints). Since
  all $n$-dimensional inner product spaces are isometric, every tangent space
  $T_pM$ is isometric to $(\mathbb{R}^n,\text{can})$, so all Riemannian manifolds of a
  given dimension share the same infinitesimal structure.
  \source{petersen:page-0021}
\end{definition}

\begin{example}[Example 1.1.1 — Euclidean space]
  \label{ex:pet-ch1-euclidean-space}
  \lean{PetersenLib.euclideanMetric}
  \leanok
  \uses{def:pet-ch1-riemannian-manifold}
  The simplest and most fundamental Riemannian manifold is Euclidean space
  $(\mathbb{R}^n,g_{\mathbb{R}^n})$. The canonical Riemannian structure $g_{\mathbb{R}^n}$ is
  defined via the tangent bundle identification $\mathbb{R}^n\times\mathbb{R}^n\simeq
  T\mathbb{R}^n$, $(p,v)\mapsto\frac{d(p+tv)}{dt}(0)$, by
  $g_{\mathbb{R}^n}((p,v),(p,w))=v\cdot w$.
  \source{petersen:page-0021}
\end{example}

\begin{definition}[Riemannian isometry]
  \label{def:pet-ch1-riemannian-isometry}
  \lean{PetersenLib.IsRiemannianIsometry}
  \leanok
  \uses{def:pet-ch1-riemannian-manifold}
  A \emph{Riemannian isometry} between Riemannian manifolds $(M,g_M)$ and
  $(N,g_N)$ is a diffeomorphism $F:M\to N$ such that $F^*g_N=g_M$, i.e.
  $g_N(DF(v),DF(w))=g_M(v,w)$ for all $p\in M$ and $v,w\in T_pM$. In that case
  $F^{-1}$ is also a Riemannian isometry.
  \source{petersen:page-0021}
\end{definition}

\begin{example}[Example 1.1.2 — inner product spaces]
  \label{ex:pet-ch1-inner-product-space}
  \lean{PetersenLib.innerProductSpaceMetric,PetersenLib.linearIsometryEquiv_isRiemannianIsometry}
  \leanok
  \uses{def:pet-ch1-riemannian-manifold,def:pet-ch1-riemannian-isometry,ex:pet-ch1-euclidean-space}
  Any finite-dimensional real vector space $V$ with an inner product becomes a
  Riemannian manifold by declaring, as with Euclidean space,
  $g((p,v),(p,w))=v\cdot w$. If $(V,g_V)$ and $(W,g_W)$ are such Riemannian
  manifolds of the same dimension, then any linear isometry $F:V\to W$ is a
  Riemannian isometry, so $(V,g_V)$ and $(W,g_W)$ are isometric. Thus
  $(\mathbb{R}^n,g_{\mathbb{R}^n})$ is (up to isometry) the only $n$-dimensional
  Riemannian manifold of this simple type.
  \source{petersen:page-0022}
\end{example}

\begin{definition}[pullback metric, Riemannian immersion]
  \label{def:pet-ch1-pullback-metric}
  \lean{PetersenLib.pullbackMetric,PetersenLib.IsRiemannianImmersion}
  \leanok
  \uses{def:pet-ch1-riemannian-manifold}
  Let $F:M\to N$ be an immersion (or embedding), with $(N,g_N)$ Riemannian. The
  \emph{pullback metric} $g_M=F^*g_N$ on $M$ is $g_M(v,w)=g_N(DF(v),DF(w))$; it
  is an inner product because $DF(v)=0$ only for $v=0$. A \emph{Riemannian
  immersion} (or \emph{Riemannian embedding}) is an immersion (or embedding)
  $F:M\to N$ with $g_M=F^*g_N$; such $F$ are also called \emph{isometric
  immersions}, though — as the following remark shows — they are almost never
  distance preserving.
  \source{petersen:page-0022}
\end{definition}

\begin{example}[Example 1.1.3 — Euclidean sphere]
  \label{ex:pet-ch1-sphere}
  \lean{PetersenLib.sphereMetric}
  \leanok
  \uses{def:pet-ch1-pullback-metric}
  The Euclidean sphere of radius $R$ is $S^n(R)=\{x\in\mathbb{R}^{n+1}\mid|x|=R\}$. The
  metric induced from the embedding $S^n(R)\hookrightarrow\mathbb{R}^{n+1}$ (via
  \cref{def:pet-ch1-pullback-metric}) is the \emph{canonical metric} on $S^n(R)$.
  The \emph{unit sphere}, or \emph{standard sphere}, is $S^n=S^n(1)\subset\mathbb{R}^{n+1}$
  with the induced metric.
  \source{petersen:page-0022}
\end{example}

\begin{remark}[nonlinear isometric immersions]
  \label{rem:pet-ch1-nonlinear-isometric-immersion}
  \lean{PetersenLib.curveIsometricEmbedding}
  \leanok
  \uses{def:pet-ch1-pullback-metric,ex:pet-ch1-euclidean-space}
  If $k<n$, linear isometric immersions $(\mathbb{R}^k,g_{\mathbb{R}^k})\to(\mathbb{R}^n,g_{\mathbb{R}^n})$
  exist but are not the only ones. Any unit-speed curve $c:\mathbb{R}\to\mathbb{R}^2$
  (i.e.\ $|\dot c(t)|=1$ for all $t$) is itself an isometric immersion; e.g.\
  $t\mapsto(\cos t,\sin t)$ is an isometric immersion and
  $t\mapsto(\log(t+\sqrt{1+t^2}),\sqrt{1+t^2})$ an isometric embedding. Given such
  a $c$, the map $F(x^1,\dots,x^k)=(c(x^1),x^2,\dots,x^k):\mathbb{R}^k\to\mathbb{R}^{k+1}$
  is a nonlinear isometric immersion (or embedding), illustrating that a
  Riemannian immersion need not be distance preserving.
  \source{petersen:page-0023}
\end{remark}

\begin{definition}[Riemannian submersion]
  \label{def:pet-ch1-riemannian-submersion}
  \lean{PetersenLib.IsRiemannianSubmersion}
  \leanok
  \uses{def:pet-ch1-riemannian-manifold}
  A \emph{Riemannian submersion} $F:(M,g_M)\to(N,g_N)$ is a submersion such that
  for each $p\in M$, $DF:\ker(DF)^\perp\to T_{F(p)}N$ is a linear isometry:
  whenever $v,w\in T_pM$ are perpendicular to $\ker(DF:T_pM\to T_{F(p)}N)$, then
  $g_M(v,w)=g_N(DF(v),DF(w))$. Equivalently, the adjoint
  $(DF_p)^*:T_{F(p)}N\to T_pM$ preserves inner products.
  \source{petersen:page-0023,petersen:page-0024}
\end{definition}

\begin{example}[Example 1.1.4 — orthogonal projections]
  \label{ex:pet-ch1-orthogonal-projection}
  \lean{PetersenLib.orthogonalProjectionSubmersion}
  \leanok
  \uses{def:pet-ch1-riemannian-submersion}
  Orthogonal projections $(\mathbb{R}^n,g_{\mathbb{R}^n})\to(\mathbb{R}^k,g_{\mathbb{R}^k})$, $k<n$, are
  Riemannian submersions.
  \source{petersen:page-0023}
\end{example}

\begin{example}[Example 1.1.5 — the Hopf fibration]
  \label{ex:pet-ch1-hopf-fibration}
  \lean{PetersenLib.hopfMap}
  \leanok
  \uses{def:pet-ch1-riemannian-submersion,ex:pet-ch1-sphere}
  The \emph{Hopf fibration} $S^3(1)\to S^2(1/2)$, viewing $S^3(1)\subset\mathbb{C}^2$
  and $S^2(1/2)\subset\mathbb{R}\oplus\mathbb{C}$, is given (F.\ Wilhelm) by
  \[
    H(z,w)=\Big(\tfrac12(|w|^2-|z|^2),\,z\bar w\Big).
  \]
  It is a Riemannian submersion.
  \source{petersen:page-0023,petersen:page-0024}
\end{example}
\begin{proof}
  \leanok
  \uses{ex:pet-ch1-hopf-fibration}
  The fibre through $(z,w)$ is $\{(e^{i\theta}z,e^{i\theta}w)\}$, so $i(z,w)$ is
  tangent to the fibre and the vectors $\lambda(-\bar w,\bar z)$, $\lambda\in\mathbb{C}$,
  span the orthogonal complement. Since $H$ extends to $\mathbb{C}^2\to\mathbb{R}\oplus\mathbb{C}$,
  restrict $H$ to the plane $(z,w)+\lambda(-\bar w,\bar z)$:
  \[
    H(z-\lambda\bar w,\,w+\lambda\bar z)
    =\Big(\tfrac12(|w+\lambda\bar z|^2-|z-\lambda\bar w|^2),\,
    (z-\lambda\bar w)(w+\lambda\bar z)\Big).
  \]
  Expanding in $\lambda,\bar\lambda$ and isolating the first-order terms gives
  \[
    DH|_{(z,w)}(\lambda(-\bar w,\bar z))=\big(2\operatorname{Re}(\bar\lambda zw),\,-\lambda\bar w^2+\bar\lambda z^2\big),
  \]
  which has the same length $|\lambda|$ as $\lambda(-\bar w,\bar z)$. Hence $DH$
  restricted to the orthogonal complement of the fibre is a linear isometry, and
  $H$ is a Riemannian submersion by \cref{def:pet-ch1-riemannian-submersion}.
\end{proof}

\begin{definition}[semi-/pseudo-Riemannian manifold; index]
  \label{def:pet-ch1-pseudo-riemannian}
  \lean{PetersenLib.PseudoRiemannianMetric,PetersenLib.pseudoRiemannianIndex}
  \leanok
  \uses{def:pet-ch1-riemannian-manifold}
  A \emph{semi-} or \emph{pseudo-Riemannian} manifold is a manifold with a
  smoothly varying symmetric bilinear form $g$ on each tangent space, assumed
  \emph{nondegenerate}: for every nonzero $v\in T_pM$ there is $w\in T_pM$ with
  $g(v,w)\ne0$. This generalizes a Riemannian metric, whose nondegeneracy follows
  from $g(v,v)>0$ for $v\ne0$. Each tangent space splits $T_pM=P\oplus N$ with $g$
  positive definite on $P$ and negative definite on $N$; the subspaces are not
  unique but their dimensions are well defined, and by continuity of $g$ nearby
  tangent spaces admit splittings of the same dimensions. The \emph{index} of a
  connected pseudo-Riemannian manifold is $\dim N$.
  \source{petersen:page-0024,petersen:page-0025}
\end{definition}

\begin{example}[Example 1.1.6 — Minkowski space]
  \label{ex:pet-ch1-minkowski}
  \lean{PetersenLib.minkowskiMetric}
  \leanok
  \uses{def:pet-ch1-pseudo-riemannian}
  For $n=n_1+n_2$ write $\mathbb{R}^{n_1,n_2}=\mathbb{R}^{n_1}\times\mathbb{R}^{n_2}$ and split vectors
  as $v=v_1+v_2$. The pseudo-Riemannian metric of index $n_2$
  $g((p,v),(p,w))=v_1\cdot w_1-v_2\cdot w_2$ is natural on $\mathbb{R}^{n_1,n_2}$. When
  $n_1=1$ or $n_2=1$ this is (a version of) Minkowski space, describing the
  geometry of special-relativistic space-time.
  \source{petersen:page-0025}
\end{example}

\begin{example}[Example 1.1.7 — hyperbolic space]
  \label{ex:pet-ch1-hyperbolic-space}
  \lean{PetersenLib.hyperbolicSpace}
  \leanok
  \uses{def:pet-ch1-pseudo-riemannian,ex:pet-ch1-minkowski,def:pet-ch1-pullback-metric}
  The \emph{hyperbolic space} $H^n(R)\subset\mathbb{R}^{n,1}$ is the branch with
  $x^{n+1}>0$ of the hyperboloid
  $(x^1)^2+\dots+(x^n)^2-(x^{n+1})^2=-R^2$, with the metric induced from the
  Minkowski metric on $\mathbb{R}^{n,1}$. This induced form is positive definite, hence a
  genuine Riemannian metric on $H^n(R)$. When $R=1$ one writes $H^n$ and calls it
  \emph{hyperbolic $n$-space}.
  \source{petersen:page-0025,petersen:page-0026}
\end{example}
\begin{proof}
  \uses{ex:pet-ch1-hyperbolic-space}
  \leanok
  A tangent vector $v=(v^1,\dots,v^{n+1})\in T_pH^n(R)$, $p\in H^n(R)$, is tangent
  to the level set of $(x^1)^2+\dots+(x^n)^2-(x^{n+1})^2$, so
  $v^1p^1+\dots+v^np^n-v^{n+1}p^{n+1}=0$, giving
  $v^{n+1}=\big(\sum_iv^ip^i\big)/p^{n+1}$. Hence the induced (Minkowski) squared
  norm is
  \[
    |v|^2=\sum_{i=1}^n(v^i)^2-(v^{n+1})^2
    =\sum_{i=1}^n(v^i)^2-\Big(\frac{\sum_iv^ip^i}{p^{n+1}}\Big)^2.
  \]
  By Cauchy--Schwarz, $\big(\sum_iv^ip^i\big)^2\le\big(\sum_i(v^i)^2\big)\big(\sum_i(p^i)^2\big)$,
  and since $p\in H^n(R)$ gives $\sum_i(p^i)^2/(p^{n+1})^2=1-(R/p^{n+1})^2$, one
  obtains
  \[
    |v|^2\ge\Big(\frac{R}{p^{n+1}}\Big)^2\sum_{i=1}^n(v^i)^2\ge0,
  \]
  with equality only for $v=0$ (since $p^{n+1}\ne0$ and the Cauchy--Schwarz
  bound is strict unless $v^1,\dots,v^n$ is proportional to $p^1,\dots,p^n$, a
  case excluded unless all $v^i=0$, forcing $v^{n+1}=0$ too). Thus the induced
  form is positive definite.
\end{proof}

\section{The Volume Form}

\begin{definition}[signed volume]
  \label{def:pet-ch1-signed-volume}
  \lean{PetersenLib.signedVolume}
  \leanok
  \uses{def:pet-ch1-riemannian-manifold}
  On an $n$-dimensional inner product space $(V,g)$ with positively oriented
  orthonormal basis $e_1,\dots,e_n$, the \emph{signed volume} of $v_1,\dots,v_n\in V$
  is
  \[
    \mathrm{vol}(v_1,\dots,v_n)=\det\big[g(v_i,e_j)\big]
    =\det\big([v_1,\dots,v_n][e_1,\dots,e_n]^T\big)=\det[v_1,\dots,v_n].
  \]
  \source{petersen:page-0026}
\end{definition}

\begin{proposition}[independence of orthonormal basis]
  \label{prop:pet-ch1-signed-volume-basis-independent}
  \lean{PetersenLib.signedVolume_orthonormal_basis_invariant}
  \leanok
  \uses{def:pet-ch1-signed-volume}
  The value $\det[g(v_i,e_j)]$ in \cref{def:pet-ch1-signed-volume} is the same for
  every positively oriented orthonormal basis $f_1,\dots,f_n$ of $(V,g)$.
  \source{petersen:page-0026}
\end{proposition}
\begin{proof}
  \uses{prop:pet-ch1-signed-volume-basis-independent}
  \leanok
  Since $[e_1,\dots,e_n]^T=[f_1,\dots,f_n]^T[f_1,\dots,f_n][e_1,\dots,e_n]^T$ is a
  factorization through the (orthogonal, determinant-$1$) change-of-basis matrix
  from $f_\bullet$ to $e_\bullet$,
  \[
    \det\big[g(v_i,f_j)\big]
    =\det\big([v_1,\dots,v_n][f_1,\dots,f_n]^T\big)
    =\det\big([v_1,\dots,v_n][f_1,\dots,f_n]^T\big)\det\big([f_1,\dots,f_n][e_1,\dots,e_n]^T\big)
  \]
  \[
    =\det\big([v_1,\dots,v_n][e_1,\dots,e_n]^T\big)=\det\big[g(v_i,e_j)\big],
  \]
  using that the determinant of the change-of-basis matrix between two positively
  oriented orthonormal bases is $1$.
\end{proof}

\begin{definition}[volume form on an oriented manifold]
  \label{def:pet-ch1-volume-form-oriented}
  \lean{PetersenLib.volumeForm}
  \leanok
  \uses{prop:pet-ch1-signed-volume-basis-independent}
  On an oriented Riemannian $n$-manifold $(M,g)$ the \emph{volume form} is the
  $n$-form
  \[
    \mathrm{vol}_g(v_1,\dots,v_n)=\mathrm{vol}(v_1,\dots,v_n)=\det\big[g(v_i,e_j)\big],
  \]
  where $e_1,\dots,e_n$ is any positively oriented orthonormal basis of
  $T_pM$ (well defined by \cref{prop:pet-ch1-signed-volume-basis-independent}).
  The notation $d\,\mathrm{vol}$ is also used, although $\mathrm{vol}$ need not be exact.
  \source{petersen:page-0026,petersen:page-0027}
\end{definition}

\begin{definition}[local volume form]
  \label{def:pet-ch1-volume-form-local}
  \lean{PetersenLib.localVolumeForm}
  \leanok
  \uses{def:pet-ch1-volume-form-oriented}
  Even when $M$ is not orientable, the volume form can be defined locally: choose
  a local \emph{orthonormal frame} $E_1,\dots,E_n$ (vector fields on a common
  domain $U\subset M$ forming a basis of $T_pM$ for all $p\in U$) and declare it
  positive, then set $\mathrm{vol}(X_1,\dots,X_n)=\det[g(X_i,E_j)]$ for vector fields
  $X_1,\dots,X_n$ on $U$.
  \source{petersen:page-0027}
\end{definition}

\begin{remark}[height times base]
  \label{rem:pet-ch1-height-base}
  \lean{PetersenLib.volumeForm_projection}
  \leanok
  \uses{def:pet-ch1-volume-form-local}
  Replacing $E_i$ by a general vector $X$ in \cref{def:pet-ch1-volume-form-local}
  gives $\mathrm{vol}(E_1,\dots,X,\dots,E_n)=g(X,E_i)$, the projection of $X$ onto
  $E_i$: the "height in the $i$th coordinate direction" times the unit base.
  \source{petersen:page-0027}
\end{remark}

\begin{remark}[integration on Riemannian manifolds]
  \label{rem:pet-ch1-integration-riemannian}
  \lean{PetersenLib.integrateOnRiemannianManifold}
  \leanok
  \uses{def:pet-ch1-volume-form-local}
  On an oriented manifold one integrates $n$-forms, so on an oriented Riemannian
  manifold one integrates a function $f$ by integrating $f\cdot\mathrm{vol}$. Every
  manifold has an open dense $O\subset M$ on which $TO\cong O\times\mathbb{R}^n$ is
  trivial, hence orientable, carrying an orthonormal frame; integrating over $O$
  therefore lets one integrate functions on any Riemannian manifold, orientable or
  not.
  \source{petersen:page-0027}
\end{remark}

\section{Groups and Riemannian Manifolds}

\subsection{Isometry Groups}

\begin{definition}[isometry group, isotropy group, homogeneous]
  \label{def:pet-ch1-isometry-group}
  \lean{PetersenLib.IsometryGroup,PetersenLib.IsotropyGroup,PetersenLib.IsHomogeneous}
  \leanok
  \uses{def:pet-ch1-riemannian-isometry}
  For a Riemannian manifold $(M,g)$, $\mathrm{Iso}(M,g)$ (or $\mathrm{Iso}(M)$)
  denotes the group of Riemannian isometries $F:(M,g)\to(M,g)$, and
  $\mathrm{Iso}_p(M,g)$ the \emph{isotropy} (or \emph{stabilizer}) subgroup at
  $p$: those $F\in\mathrm{Iso}(M,g)$ with $F(p)=p$. $(M,g)$ is \emph{homogeneous}
  if $\mathrm{Iso}(M,g)$ acts transitively: for all $p,q\in M$ there is
  $F\in\mathrm{Iso}(M,g)$ with $F(p)=q$.
  \source{petersen:page-0027}
\end{definition}

\begin{example}[Example 1.3.1 — isometries of Euclidean space]
  \label{ex:pet-ch1-isometry-group-euclidean}
  \lean{PetersenLib.isometryGroup_euclideanSpace}
  \leanok
  \uses{def:pet-ch1-isometry-group,def:pet-ch1-riemannian-isometry}
  $\mathrm{Iso}(\mathbb{R}^n,g_{\mathrm{eu}})=\mathbb{R}^n\rtimes O(n)=\{F(x)=v+Ox\mid
  v\in\mathbb{R}^n,\,O\in O(n)\}$, with $v,O$ uniquely determined by $F$. The isotropy
  $\mathrm{Iso}_p$ is isomorphic to $O(n)$ for any $p$, and
  $\mathbb{R}^n\simeq\mathrm{Iso}/\mathrm{Iso}_p$; in general every homogeneous space is a
  quotient $M=\mathrm{Iso}/\mathrm{Iso}_p$.
  \source{petersen:page-0028}
\end{example}
\begin{proof}
  \leanok
  \uses{ex:pet-ch1-isometry-group-euclidean,rem:pet-ch1-ex-17,rem:pet-ch1-ex-19}
  Maps $x\mapsto v+Ox$ with $v\in\mathbb{R}^n$, $O\in O(n)$ are visibly isometries.
  Conversely, let $F\in\mathrm{Iso}(\mathbb{R}^n,g_{\mathrm{eu}})$. (Petersen
  derives the converse from the uniqueness principle for Riemannian
  isometries, Prop.\ 5.6.2 of the source, outside this chapter; we argue
  instead through arc length and the Mazur--Ulam theorem.) A Riemannian
  isometry preserves the length of curves
  (\cref{rem:pet-ch1-ex-17}). For $x,y\in\mathbb{R}^n$ the image under
  $F$ of the straight segment from $x$ to $y$ is a smooth curve from $F(x)$
  to $F(y)$ of length $|y-x|$, and every curve is at least as long as the
  distance between its endpoints (\cref{rem:pet-ch1-ex-19}), so $|F(x)-F(y)|\le|x-y|$. Applying the
  same argument to $F^{-1}$ gives the reverse inequality, so $F$ is a
  surjective distance-preserving map of $\mathbb{R}^n$. By the Mazur--Ulam
  theorem $F$ is affine, $F(x)=F(0)+O(x)$ with $O$ linear and
  distance-preserving with $O(0)=0$, hence $O\in O(n)$; take $v=F(0)$.
\end{proof}

\begin{example}[Example 1.3.2 — isometries of the sphere]
  \label{ex:pet-ch1-isometry-group-sphere}
  \lean{PetersenLib.isometryGroup_sphere}
  \leanok
  \uses{def:pet-ch1-isometry-group,ex:pet-ch1-sphere}
  $\mathrm{Iso}(S^n(R),g_{S^n(R)})=O(n+1)=\mathrm{Iso}_0(\mathbb{R}^{n+1},g_{\mathrm{eu}})$.
  The isotropy groups are isomorphic to $O(n)$ (elements of $O(n+1)$ fixing a
  line), so $S^n\simeq O(n+1)/O(n)$.
  \source{petersen:page-0028}
\end{example}
\begin{proof}
  \uses{ex:pet-ch1-isometry-group-sphere,rem:pet-ch1-ex-20,rem:pet-ch1-ex-17}
  \leanok
  Clearly $O(n+1)\subset\mathrm{Iso}(S^n(R),g_{S^n(R)})$. Conversely, let
  $F\in\mathrm{Iso}(S^n(R),g_{S^n(R)})$. Petersen's text assembles the matrix
  $O=\big[\tfrac1RF(Re_1)\ DF|_{e_1}(e_2)\ \cdots\ DF|_{e_1}(e_{n+1})\big]$ and
  invokes the uniqueness principle (Prop.\ 5.6.2 of the source) to conclude
  $F=O$; since that principle lies outside Chapter~1, we argue instead through
  the length-minimization characterization of great circles
  (Exercise~1.6.20). For $x,y\in S^n(R)$ with angle
  $\theta=\arccos(\langle x,y\rangle/R^2)\in(0,\pi)$, the great circle from
  $x$ to $y$ has induced length $R\theta$; since $F$ preserves induced arc
  length (Exercise~1.6.17 and $F^*g=g$), its image is a curve from $F(x)$ to
  $F(y)$ of the same length, and the lower bound of Exercise~1.6.20(4) gives
  $R\,\arccos(\langle F(x),F(y)\rangle/R^2)\le R\theta$. Hence
  $\langle x,y\rangle\le\langle F(x),F(y)\rangle$ (the degenerate cases
  $\langle x,y\rangle=\pm R^2$ follow from Cauchy--Schwarz); applying the
  same to $F^{-1}$ yields $\langle F(x),F(y)\rangle=\langle x,y\rangle$.
  Now let $(e_i)$ be an orthonormal basis and define $O$ linearly by
  $O(e_i)=\tfrac1R F(Re_i)$. By the preservation of inner products the family
  $(O(e_i))$ is orthonormal, so $O\in O(n+1)$; and for every $x\in S^n(R)$,
  $\langle O(e_i),F(x)\rangle=\tfrac1R\langle F(Re_i),F(x)\rangle
  =\tfrac1R\langle Re_i,x\rangle=\langle O(e_i),O(x)\rangle$ for all $i$,
  so $F(x)=O(x)$.
\end{proof}

\begin{example}[Example 1.3.3 — isometries of hyperbolic space]
  \label{ex:pet-ch1-isometry-group-hyperbolic}
  \lean{PetersenLib.isometryGroup_hyperbolicSpace}
  \leanok
  \uses{def:pet-ch1-isometry-group,ex:pet-ch1-hyperbolic-space}
  Let $O(n,1)=\{L:\mathbb{R}^{n,1}\to\mathbb{R}^{n,1}\mid g(Lv,Lv)=g(v,v)\}$ and
  $O^+(n,1)\subset O(n,1)$ those $L$ preserving $x^{n+1}>0$ on timelike
  vectors. A permutation $F$ of $H^n(R)$ is a Riemannian isometry of
  $(H^n(R),g_{H^n(R)})$ if and only if it is the restriction to $H^n(R)$ of
  an orthochronous transformation $L\in O^+(n,1)$, so $\mathrm{Iso}(H^n(R))$
  is realized inside $O^+(n,1)$; the isotropy of $Re_{n+1}$ is identified with
  $O(n)$, and $O^+(n,1)$ acts transitively on $H^n(R)$.
  \source{petersen:page-0028,petersen:page-0029}
\end{example}
\begin{proof}
  \uses{ex:pet-ch1-isometry-group-hyperbolic,rem:pet-ch1-ex-21,rem:pet-ch1-ex-17}
  \leanok
  If $L\in O^+(n,1)$, then $L$ preserves the defining equation
  $g(x,x)=-R^2$ and the sheet condition $x^{n+1}>0$, so it restricts to a
  bijection of $H^n(R)$, smooth in the global chart; since the metric of
  $H^n(R)$ is induced from the Minkowski form, which (by polarization) $L$
  preserves as a bilinear form, the restriction is a Riemannian isometry.

  Conversely let $F\in\mathrm{Iso}(H^n(R))$. Petersen's text assembles the
  matrix with columns
  $DF|_{e_{n+1}}(e_1),\dots,DF|_{e_{n+1}}(e_n),\tfrac1RF(Re_{n+1})$ and
  invokes the uniqueness principle (Prop.~5.6.2 of the source) to conclude
  that $F$ agrees with it; since that principle lies outside Chapter~1, we
  argue instead through the length-minimization characterization of
  hyperbolas (Exercise~1.6.21). For $x,y\in H^n(R)$ with
  $\theta=\operatorname{arcosh}(-g(x,y)/R^2)>0$ the hyperbola
  $s\mapsto x\cosh s+Rv\sinh s$, with $v$ the unit tangent at $x$ toward
  $y$, reaches $y$ at $s=\theta$ and has induced length $R\theta$; since
  $F$ preserves induced arc length (Exercise~1.6.17 and $F^*g=g$), its
  image is a curve from $F(x)$ to $F(y)$ of the same length, and the lower
  bound of Exercise~1.6.21(4) gives
  $R\operatorname{arcosh}(-g(F(x),F(y))/R^2)\le R\theta$; as both pairings
  satisfy $-g(\cdot,\cdot)/R^2\ge1$ and $\operatorname{arcosh}$ is monotone,
  $-g(F(x),F(y))\le-g(x,y)$. Applying the same to $F^{-1}$ yields
  $g(F(x),F(y))=g(x,y)$ for all $x,y\in H^n(R)$.

  Now the $n+1$ points $W_0=(0,R)$ and $W_i=(R\sinh(1)\,e_i,R\cosh(1))$,
  $i=1,\dots,n$, of $H^n(R)$ form a basis of $\mathbb{R}^{n,1}$: in a
  vanishing linear combination the first components force the coefficients
  of the $W_i$ (as $R\sinh1\neq0$ and the $e_i$ are independent), and the
  last coordinate then forces that of $W_0$. Define $L$ linearly by
  $L(W_k)=F(W_k)$. Since $F$ preserves pairings of hyperboloid points,
  $g(Lv,Lw)=g(v,w)$ holds on basis pairs, hence for all $v,w$ by
  bilinearity, so $L\in O(n,1)$. Moreover $L$ is injective — if $Lv=0$
  then $g(v,\cdot)\equiv0$, while the Minkowski form is nondegenerate
  ($g(v,(v_1,-v^{n+1}))=|v_1|^2+(v^{n+1})^2>0$ for $v\neq0$) — hence
  bijective. For any $x\in H^n(R)$ and each $k$,
  $g(F(x),LW_k)=g(F(x),F(W_k))=g(x,W_k)=g(Lx,LW_k)$; as $(LW_k)$ is again
  a basis, nondegeneracy gives $F(x)=L(x)$, so $F$ is the restriction of
  $L$. Finally $L$ is orthochronous: an upper timelike vector $v$ is a
  positive multiple $v=c\,p$ of a point $p\in H^n(R)$ (with
  $c=\sqrt{-g(v,v)}/R>0$), and $(Lv)^{n+1}=c\,(F(p))^{n+1}>0$.
\end{proof}

\subsection{Lie Groups}

\begin{definition}[left-invariant metric]
  \label{def:pet-ch1-left-invariant-metric}
  \lean{PetersenLib.leftInvariantMetric}
  \leanok
  \uses{def:pet-ch1-riemannian-isometry}
  For a Lie group $G$, left translation trivializes $TG\simeq G\times T_eG$, so
  any inner product on $T_eG$ induces a Riemannian metric on $G$ for which every
  left translation $L_x$ is a Riemannian isometry — a \emph{left-invariant}
  metric. Every Riemannian metric on $G$ for which all left translations are
  isometries arises this way. Distinct inner products on $T_eG$ need not give
  isometric metrics, so a Lie group need not carry a canonical metric.
  \source{petersen:page-0029}
\end{definition}

\begin{remark}[Riemannian submersions onto homogeneous quotients]
  \label{rem:pet-ch1-quotient-submersion}
  \lean{PetersenLib.homogeneousQuotientMetric}
  \leanok
  \uses{def:pet-ch1-riemannian-submersion,def:pet-ch1-left-invariant-metric}
  If $H\subset G$ is a compact subgroup, $G/H=\{gH\}$ is a manifold. Given a
  Riemannian metric on $G$ for which right translations by elements of $H$ act by
  isometries, there is a unique Riemannian metric on $G/H$ making $G\to G/H$ a
  Riemannian submersion; if the metric on $G$ is also left-invariant, $G$ acts by
  isometries on $G/H$, making $G/H$ homogeneous.
  \source{petersen:page-0029}
\end{remark}

\begin{example}[Example 1.3.4 — Fubini--Study metric]
  \label{ex:pet-ch1-fubini-study}
  \lean{PetersenLib.fubiniStudyMetric,PetersenLib.ComplexProjectiveSpace,PetersenLib.projSphere,PetersenLib.chartCP,PetersenLib.fubiniStudyMetricComplexProjectiveSpace}
  \leanok
  \uses{rem:pet-ch1-quotient-submersion,ex:pet-ch1-sphere}
  The circle $S^1=\{\lambda\in\mathbb{C}\mid|\lambda|=1\}$ acts by isometric complex
  scalar multiplication on $S^{2n+1}(1)\subset\mathbb{C}^{n+1}$, with quotient
  $S^{2n+1}/S^1=\mathbb{CP}^n$; the induced metric on $\mathbb{CP}^n$ making
  $S^{2n+1}\to\mathbb{CP}^n$ a Riemannian submersion is the \emph{Fubini--Study metric}.
  When $n=1$ this is the Hopf fibration $S^3(1)\to\mathbb{CP}^1=S^2(1/2)$.
  \source{petersen:page-0029}
\end{example}

\begin{example}[Example 1.3.5 — $SU(2)$ and the Berger spheres]
  \label{ex:pet-ch1-su2-berger}
  \lean{PetersenLib.suTwoMetric,PetersenLib.bergerSphereMetric}
  \leanok
  \uses{def:pet-ch1-left-invariant-metric,ex:pet-ch1-sphere}
  $SU(2)=\{A\in M_{2\times2}(\mathbb{C})\mid\det A=1,A^*=A^{-1}\}
  =\{\begin{bsmallmatrix}z&w\\-\bar w&\bar z\end{bsmallmatrix}:|z|^2+|w|^2=1\}=S^3(1)$,
  with Lie algebra $\mathfrak{su}(2)$ spanned by
  $X_1=\mathrm{diag}(i,-i)$, $X_2=\begin{bsmallmatrix}0&1\\-1&0\end{bsmallmatrix}$,
  $X_3=\begin{bsmallmatrix}0&i\\i&0\end{bsmallmatrix}$, viewed as left-invariant
  vector fields. Declaring $X_1,X_2,X_3$ orthonormal gives a left-invariant
  metric on $SU(2)$ isometric to $S^3(1)$. Declaring them orthogonal, with
  $|X_1|=\varepsilon$ and $X_2,X_3$ unit, gives the \emph{Berger spheres}
  $g_\varepsilon$ on $SU(2)=S^3$; scalar multiplication on $S^3\subset\mathbb{C}^2$
  corresponds to left multiplication by $\mathrm{diag}(e^{i\theta},e^{-i\theta})$,
  so $X_1$ is tangent to the Hopf circle action, and $g_\varepsilon$ is obtained
  from the canonical metric by scaling along the Hopf fibre by $\varepsilon^2$.
  \source{petersen:page-0029,petersen:page-0030}
\end{example}

\subsection{Covering Maps}

\begin{remark}[covering-induced and quotient metrics]
  \label{rem:pet-ch1-covering-metric}
  \lean{PetersenLib.coveringInducedMetric,PetersenLib.quotientMetric}
  \leanok
  \uses{def:pet-ch1-riemannian-submersion,def:pet-ch1-pullback-metric}
  A covering map $F:M\to N$ is both an immersion and a submersion, so a
  Riemannian metric on $N$ pulls back to $M$, making $F$ an isometric immersion
  (\emph{Riemannian covering}); since $\dim M=\dim N$, $F$ is in fact a local
  isometry. If $F$ is a normal covering with deck group $\Gamma$, then $\Gamma$
  acts by isometries on the pullback metric. Conversely, if $M$ is Riemannian and
  $\Gamma$ acts on $M$ by isometries with $N=M/\Gamma$, there is a unique
  Riemannian metric on $N$ making the quotient map a local isometry.
  \source{petersen:page-0030,petersen:page-0031}
\end{remark}

\begin{example}[Example 1.3.6 — the flat torus]
  \label{ex:pet-ch1-flat-torus}
  \lean{PetersenLib.flatTorus}
  \leanok
  \uses{rem:pet-ch1-covering-metric}
  Fixing a basis $v_1,v_2$ of $\mathbb{R}^2$, $\mathbb{Z}^2$ acts by isometric translations
  $(n,m)\mapsto(x\mapsto x+nv_1+mv_2)$; the quotient is a torus $T^2$ with a flat
  metric depending on $|v_1|,|v_2|,\angle(v_1,v_2)$, so distinct choices need not
  be isometric. $T^2$ is itself an abelian Lie group and these metrics are
  bi-invariant.
  \source{petersen:page-0031}
\end{example}

\begin{example}[Example 1.3.7 — real projective space]
  \label{ex:pet-ch1-real-projective-space}
  \lean{PetersenLib.realProjectiveSpaceCovering}
  \leanok
  \uses{rem:pet-ch1-covering-metric,ex:pet-ch1-sphere}
  The involution $-I$ on $S^n(1)\subset\mathbb{R}^{n+1}$ is an isometry, inducing a
  Riemannian covering $S^n\to\mathbb{RP}^n$.
  \source{petersen:page-0031}
\end{example}

\section{Local Representations of Metrics}

\subsection{Einstein Summation Convention}

\begin{convention}[Einstein summation]
  \label{conv:pet-ch1-einstein-summation}
  \lean{PetersenLib.einsteinSummationConvention}
  \leanok
  For a vector space $V$ (such as a tangent space) basis vectors carry
  subscripts, $e_1,\dots,e_n$, and coefficients carry superscripts, so
  $v=\sum_iv^ie_i=v^ie_i$ (repeated sub/superscript pairs are summed). The dual
  basis $e^i\in V^*=\mathrm{Hom}(V,\mathbb{R})$ satisfies $e^i(e_j)=\delta^i_j$, and
  $v^i=e^i(v)$; dual bases thus also carry superscripts. A linear map
  $L:V\to V$ has matrix $[L^i_j]$ determined by $L(e_i)=L^i_je_j$, i.e.
  $L^i_j=e^i(L(e_j))$: superscripts index rows, subscripts columns. On a
  manifold, coordinates $x^i$ carry superscripts, so coordinate vector fields
  $\partial_i=\partial/\partial x^i$ carry subscripts and form a basis of $T_pM$,
  while the dual $1$-forms $dx^i$, satisfying $dx^i(\partial_j)=\delta^i_j$, form
  the dual basis of $T_p^*M$. For a curve $c(t)=(x^i(t))=x^i(t)e_i$ in
  $\mathbb{R}^n$, the velocity $\dot c(t)=(\dot x^i(t))=\dot x^i(t)\partial_i$ is
  coordinate-independent because $\dot x^i(t)=dx^i(\dot c(t))$ by the chain rule,
  matching the definition of the components of a vector in a given basis.
  \source{petersen:page-0031,petersen:page-0032,petersen:page-0033}
\end{convention}

\subsection{Coordinate Representations}

\begin{definition}[coordinate representation of a metric]
  \label{def:pet-ch1-coordinate-metric}
  \lean{PetersenLib.metricCoordinateComponents}
  \leanok
  \uses{def:pet-ch1-riemannian-manifold,conv:pet-ch1-einstein-summation}
  On coordinates $x=(x^1,\dots,x^n)$ on $U\subset M$, $1$-forms multiply to
  bilinear forms via $\theta_1\cdot\theta_2(v,w)=\theta_1(v)\theta_2(w)=
  (\theta_1\otimes\theta_2)(v,w)$ (in general $\theta_1\cdot\theta_2\ne
  \theta_2\cdot\theta_1$). A Riemannian metric $g$ is then
  \[
    g=g(\partial_i,\partial_j)\,dx^i\cdot dx^j,
  \]
  since $g(v,w)=g(dx^i(v)\partial_i,dx^j(w)\partial_j)=g(\partial_i,\partial_j)\,
  dx^i(v)\,dx^j(w)$. The functions $g_{ij}:=g(\partial_i,\partial_j)=g_{ji}$ form a
  positive-definite symmetric matrix parametrized over $U$: the \emph{local
  representation} of $g$.
  \source{petersen:page-0033}
\end{definition}

\begin{example}[Example 1.4.1 — canonical metric on $\mathbb{R}^n$]
  \label{ex:pet-ch1-euclidean-coords}
  \lean{PetersenLib.euclideanMetricCoordinates}
  \leanok
  \uses{def:pet-ch1-coordinate-metric,ex:pet-ch1-euclidean-space}
  In the identity chart, $g_{\mathbb{R}^n}=\delta_{ij}\,dx^i\,dx^j=\sum_{i=1}^n(dx^i)^2$.
  \source{petersen:page-0033}
\end{example}

\begin{example}[Example 1.4.2 — polar coordinates]
  \label{ex:pet-ch1-polar-coords}
  \lean{PetersenLib.polarCoordinateMetric}
  \leanok
  \uses{def:pet-ch1-coordinate-metric,ex:pet-ch1-euclidean-coords}
  On $\mathbb{R}^2$ minus a half-line, in polar coordinates $(r,\theta)$ the canonical
  metric is $g=dr^2+r^2\,d\theta^2$, i.e.\ $g_{rr}=1$, $g_{r\theta}=g_{\theta r}=0$,
  $g_{\theta\theta}=r^2$.
  \source{petersen:page-0034}
\end{example}
\begin{proof}
  \uses{ex:pet-ch1-polar-coords}
  \leanok
  With $x=r\cos\theta$, $y=r\sin\theta$, $dx=\cos\theta\,dr-r\sin\theta\,d\theta$
  and $dy=\sin\theta\,dr+r\cos\theta\,d\theta$, so
  \[
    dx^2+dy^2=(\cos\theta\,dr-r\sin\theta\,d\theta)^2+(\sin\theta\,dr+r\cos\theta\,d\theta)^2
  \]
  \[
    =(\cos^2\theta+\sin^2\theta)\,dr^2
    +\big(r\cos\theta\sin\theta-r\cos\theta\sin\theta\big)(dr\,d\theta+d\theta\,dr)
    +r^2(\sin^2\theta+\cos^2\theta)\,d\theta^2
  \]
  \[
    =dr^2+r^2\,d\theta^2.\qedhere
  \]
\end{proof}

\subsection{Frame Representations}

\begin{definition}[frame representation of a metric]
  \label{def:pet-ch1-frame-metric}
  \lean{PetersenLib.frameRepresentation}
  \leanok
  \uses{def:pet-ch1-riemannian-manifold}
  A \emph{frame} $X_1,\dots,X_n$ on $U\subset M$ is a collection of $n=\dim M$
  linearly independent vector fields on $U$; the \emph{coframe} $\sigma^1,\dots,
  \sigma^n$ is the dual $1$-forms, $\sigma^i(X_j)=\delta^i_j$. The metric is then
  represented as $g=g_{ij}\,\sigma^i\sigma^j=g(X_i,X_j)\,\sigma^i\sigma^j$.
  \source{petersen:page-0034}
\end{definition}

\begin{example}[Example 1.4.3 — left-invariant metrics via a coframe]
  \label{ex:pet-ch1-left-invariant-coframe}
  \lean{PetersenLib.leftInvariantCoframeMetric}
  \leanok
  \uses{def:pet-ch1-frame-metric,def:pet-ch1-left-invariant-metric,ex:pet-ch1-su2-berger}
  Given a frame $X_1,\dots,X_n$ of left-invariant vector fields on a Lie group
  $G$ with dual coframe $\sigma^1,\dots,\sigma^n$, a left-invariant metric $g$
  depends only on its values on $T_eG$ and is $g=g_{ij}\sigma^i\sigma^j$ for a
  positive-definite symmetric constant matrix $(g_{ij})$. If $X_i$ form an
  orthonormal basis for $\mathfrak{g}$, $g=(\sigma^1)^2+\dots+(\sigma^n)^2$. The
  Berger sphere is $g_\varepsilon=\varepsilon^2(\sigma^1)^2+(\sigma^2)^2+(\sigma^3)^2$.
  \source{petersen:page-0034,petersen:page-0035}
\end{example}

\begin{example}[Example 1.4.4 — surfaces of revolution]
  \label{ex:pet-ch1-surface-of-revolution}
  \lean{PetersenLib.surfaceOfRevolutionMetric}
  \leanok
  \uses{def:pet-ch1-frame-metric}
  A surface of revolution has profile curve $c(t)=(r(t),0,z(t)):I\to\mathbb{R}^3$
  ($r(t)>0$), revolved about the $z$-axis:
  $f(t,\theta)=(r(t)\cos\theta,r(t)\sin\theta,z(t))$, with natural frame
  $\partial_t,\partial_\theta$ and coframe $dt,d\theta$. The induced metric from
  $dx^2+dy^2+dz^2$ is $g=(\dot r^2+\dot z^2)dt^2+r^2\,d\theta^2$; if $c$ is
  parametrized by arc length ($\dot r^2+\dot z^2=1$), $g=dt^2+r^2\,d\theta^2$.
  \source{petersen:page-0035}
\end{example}
\begin{proof}
  \uses{ex:pet-ch1-surface-of-revolution}
  \leanok
  $dx=\dot r\cos\theta\,dt-r\sin\theta\,d\theta$,
  $dy=\dot r\sin\theta\,dt+r\cos\theta\,d\theta$, $dz=\dot z\,dt$, so
  \[
    dx^2+dy^2+dz^2
    =(\dot r\cos\theta\,dt-r\sin\theta\,d\theta)^2+(\dot r\sin\theta\,dt+r\cos\theta\,d\theta)^2+(\dot z\,dt)^2
  \]
  \[
    =(\dot r^2+\dot z^2)\,dt^2+r^2\,d\theta^2,
  \]
  by the same cross-term cancellation as in \cref{ex:pet-ch1-polar-coords}.
\end{proof}

\begin{remark}[rotationally symmetric metrics]
  \label{rem:pet-ch1-rotationally-symmetric}
  \lean{PetersenLib.rotationallySymmetricMetric}
  \leanok
  \uses{ex:pet-ch1-surface-of-revolution}
  On the abstract manifold $I\times S^1$, with frame $\partial_t,\partial_\theta$,
  one defines \emph{rotationally symmetric} metrics $g=\eta^2(t)dt^2+\rho^2(t)d\theta^2$
  ($\eta,\rho$ independent of $\theta$); a change of coordinates on $I$ generally
  reduces to $\eta\equiv1$. Not every rotationally symmetric metric comes from a
  surface of revolution: if $dt^2+r^2d\theta^2$ is such a surface then
  $\dot z^2+\dot r^2=1$, so $|\dot r|\le1$.
  \source{petersen:page-0036}
\end{remark}

\begin{example}[Example 1.4.5 — the round sphere as a surface of revolution]
  \label{ex:pet-ch1-sphere-revolution}
  \lean{PetersenLib.sphereAsSurfaceOfRevolution}
  \leanok
  \uses{rem:pet-ch1-rotationally-symmetric,ex:pet-ch1-sphere,ex:pet-ch1-hyperbolic-space}
  $S^2(R)\subset\mathbb{R}^3$ is the surface of revolution of
  $t\mapsto R(\sin(t/R),0,\cos(t/R))$, with metric
  $dt^2+R^2\sin^2(t/R)\,d\theta^2$; since $R\sin(t/R)\to t$ as $R\to\infty$,
  large spheres approximate Euclidean space. Formally replacing $R$ by $iR$
  yields $dt^2+R^2\sinh^2(t/R)\,d\theta^2$, which coincides with the induced
  metric on $H^2(R)$ (\cref{ex:pet-ch1-hyperbolic-space}) obtained by revolving
  $t\mapsto R(\sinh(t/R),0,\cosh(t/R))$ around the $z$-axis in $\mathbb{R}^{2,1}$.
  \source{petersen:page-0036}
\end{example}

\begin{definition}[the functions $\mathrm{sn}_k$, $\mathrm{cs}_k$]
  \label{def:pet-ch1-snk-csk}
  \lean{PetersenLib.snFunction,PetersenLib.csFunction}
  \leanok
  \uses{rem:pet-ch1-rotationally-symmetric}
  Let $\mathrm{sn}_k(t)$ be the unique solution of $\ddot x+kx=0$, $x(0)=0$,
  $\dot x(0)=1$, and $\mathrm{cs}_k(t)$ the unique solution of $\ddot x+kx=0$,
  $x(0)=1$, $\dot x(0)=0$. They satisfy
  $\dot{\mathrm{sn}}_k=\mathrm{cs}_k$, $\dot{\mathrm{cs}}_k=-k\,\mathrm{sn}_k$,
  $\mathrm{cs}_k^2+k\,\mathrm{sn}_k^2=1$. The family
  $dt^2+\mathrm{sn}_k^2(t)\,d\theta^2$ is $\mathbb{R}^2$ when $k=0$, $S^2(1/\sqrt k)$ when
  $k>0$, and $H^2(1/\sqrt{-k})$ when $k<0$.
  \source{petersen:page-0036,petersen:page-0037}
\end{definition}

\subsection{Polar Versus Cartesian Coordinates}

\begin{definition}[higher-dimensional rotationally symmetric metric]
  \label{def:pet-ch1-rot-sym-general}
  \lean{PetersenLib.rotationallySymmetricMetricHighDim}
  \leanok
  \uses{rem:pet-ch1-rotationally-symmetric}
  On $I\times S^{n-1}$, with $ds^2_{n-1}$ the canonical metric on
  $S^{n-1}(1)\subset\mathbb{R}^n$, one considers \emph{rotationally symmetric metrics}
  $dt^2+\rho^2(t)ds^2_{n-1}$ (a special class of \emph{warped products}), with
  $\rho(0)=0$, $\rho(t)>0$ for $t>0$; the question is when such a metric extends
  smoothly across the origin of $\mathbb{R}^n$.
  \source{petersen:page-0037}
\end{definition}

\begin{lemma}[polar-coordinate change formula]
  \label{lem:pet-ch1-polar-coordinate-change}
  \lean{PetersenLib.polarToCartesianFormula}
  \leanok
  \uses{def:pet-ch1-rot-sym-general}
  Write $x=ts$, $x\in\mathbb{R}^n$, $t=|x|>0$, $s\in S^{n-1}(1)$, so
  $\sum(s^i)^2=1$ and $\sum s^ids^i=0$. Then $\sum(dx^i)^2=dt^2+t^2ds_{n-1}^2$,
  hence $dt=\tfrac1t\sum x^idx^i$ and $ds_{n-1}^2=\big(\sum(dx^i)^2-dt^2\big)/t^2$,
  giving
  \[
    dt^2+\rho^2(t)ds_{n-1}^2
    =\Big(\frac1{t^2}-\frac{\rho^2(t)}{t^4}\Big)\Big(\sum s^idx^i\Big)^2
    +\frac{\rho^2(t)}{t^2}\sum(dx^i)^2.
  \]
  \source{petersen:page-0037,petersen:page-0038}
\end{lemma}
\begin{proof}
  \uses{lem:pet-ch1-polar-coordinate-change}
  \leanok
  From $x^i=ts^i$, $dx^i=s^idt+tds^i$, so
  $\sum(dx^i)^2=\sum(s^i)^2dt^2+t^2\sum(ds^i)^2+2t\,dt\sum s^ids^i
  =dt^2+t^2ds_{n-1}^2$, using $\sum(s^i)^2=1$ and $\sum s^ids^i=0$. From
  $t^2=\sum(x^i)^2$, $2t\,dt=2\sum x^idx^i$, giving $dt=\tfrac1t\sum x^idx^i$.
  Solving $\sum(dx^i)^2=dt^2+t^2ds_{n-1}^2$ for $ds_{n-1}^2$ gives the stated
  formula, and substituting into $dt^2+\rho^2(t)ds_{n-1}^2$ and regrouping the
  $dt^2=\tfrac1{t^2}(\sum s^idx^i)^2$ term against the $\sum(dx^i)^2$ term yields
  the claimed identity.
\end{proof}

\begin{theorem}[smoothness at the origin]
  \label{thm:pet-ch1-smoothness-criterion}
  \lean{PetersenLib.rotationallySymmetricSmoothnessCriterion}
  \leanok
  \uses{lem:pet-ch1-polar-coordinate-change}
  In the notation of \cref{lem:pet-ch1-polar-coordinate-change}, with ambient
  dimension $n\ge1$, the metric
  $dt^2+\rho^2(t)ds_{n-1}^2$ extends to a smooth metric on $\mathbb{R}^n$ near the origin
  $t=0$ if and only if $\rho(0)=0$, $\dot\rho(0)=1$, and all even-order
  derivatives of $\rho$ vanish at $0$: $\rho^{(2l)}(0)=0$ for $l\ge1$.
  \source{petersen:page-0038,petersen:page-0039}
\end{theorem}
\begin{proof}
  \leanok
  \uses{thm:pet-ch1-smoothness-criterion}
  By \cref{lem:pet-ch1-polar-coordinate-change} it suffices that
  $F_1=\rho^2(t)/t^2$ and $F_2=1/t^2-\rho^2(t)/t^4$ extend smoothly in $x$
  across the origin with $F_1(0)>0$. (Petersen argues by Taylor expansion
  that both coefficients are approximable to any order by polynomials in
  $t^2$; approximability alone does not give smoothness without uniform
  control of the remainders, so we argue instead by Hadamard quotients and
  Whitney's even-function theorem, following the same idea.)

  Assume $\rho(0)=0$, $\dot\rho(0)=1$, $\rho^{(2l)}(0)=0$ for $l\ge1$. Split
  $\rho=\rho_o+\rho_e$ into odd and even parts; the hypotheses say exactly
  that every derivative of $\rho_e$ vanishes at $0$, i.e.\ $\rho_e$ is flat
  at $0$. The Hadamard quotient $f\mapsto\big(t\mapsto\int_0^1\dot
  f(ts)\,ds\big)$ preserves smoothness and flatness and exchanges parity,
  and satisfies $f(t)=f(0)+t\int_0^1\dot f(ts)\,ds$. Applying it repeatedly
  produces smooth functions $K$ (even, $K(0)=\dot\rho(0)=1$), $M$ (flat),
  $H$ (even) and $X$ (flat) with
  \[
    \rho(t)=t\,(K(t)+M(t)),\qquad 1-K^2(t)=t^2H(t),\qquad M(t)=t^2X(t),
  \]
  whence for $t\ne0$
  \[
    F_1=(K+M)^2,\qquad F_2=H-X\,(2K+M).
  \]
  The right-hand sides are globally defined smooth functions of $t$ of the
  form (even)${}+{}$(flat). By Whitney's even-function theorem a smooth even
  function of $t$ is a smooth function of $t^2$, and a smooth flat function
  of $t$ is one as well; since $t^2=\sum(x^i)^2$ is smooth in $x$, both
  coefficients extend smoothly across the origin, with
  $F_1(0)=K(0)^2=1>0$.

  Conversely, suppose $F_1,F_2$ extend smoothly. From
  $F_2=(1-F_1)/t^2$ we get $F_1=1-t^2F_2$, so $F_1(0)=1$ by continuity.
  Restricting along a ray through a unit vector $e$, the function
  $a(t)=F_1(te)$ is smooth near $0$, even (both signs of $t$ give
  $\rho(|t|)^2/t^2$), and $a(0)=1$. Correcting $a$ outside a small interval
  by an even cutoff makes it globally smooth, even and positive, so
  $r(t)=t\sqrt{a(t)}$ is a globally smooth \emph{odd} function, and for
  small $t>0$, $r(t)=t\sqrt{\rho(t)^2/t^2}=\rho(t)$ by positivity of
  $\rho$. Two smooth functions agreeing on $(0,\delta)$ have equal
  derivatives at $0$: hence $\rho(0)=r(0)=0$,
  $\dot\rho(0)=\dot r(0)=\sqrt{a(0)}=1$, and every even-order derivative of
  $\rho$ at $0$ vanishes because $r$ is odd.
\end{proof}

\begin{example}[Example 1.4.6 — space forms]
  \label{ex:pet-ch1-space-forms}
  \lean{PetersenLib.spaceForm,PetersenLib.spaceFormIsometry,PetersenLib.spaceFormIsometry_neg}
  \leanok
  \uses{thm:pet-ch1-smoothness-criterion,def:pet-ch1-snk-csk,ex:pet-ch1-sphere,ex:pet-ch1-hyperbolic-space}
  The hypotheses of \cref{thm:pet-ch1-smoothness-criterion} hold for all metrics
  $dt^2+\mathrm{sn}_k^2(t)ds_{n-1}^2$, $t\in[0,\infty)$ if $k\le0$ and
  $t\in[0,\pi/\sqrt k]$ if $k>0$. The resulting Riemannian manifolds, the
  \emph{space forms} $S^n_k$ of dimension $n$ and curvature $k$, are isometric to
  $\mathbb{R}^n$ ($k=0$), $S^n(R)$ ($k=1/R^2$), or $H^n(R)$ ($k=-1/R^2$).
  \source{petersen:page-0039,petersen:page-0040}
\end{example}
\begin{proof}
  \leanok
  \uses{ex:pet-ch1-space-forms}
  For $k=1/R^2$: the map $G:S^{n-1}\times(0,R\pi)\to\mathbb{R}^{n+1}$,
  $G(s,r)=R(s\sin(r/R),\cos(r/R))$, has image in $S^n(R)$. Writing
  $x=Rs\sin(r/R)$, $t=R\cos(r/R)$, and using $\sum(s^i)^2=1$,
  $\sum s^ids^i=0$,
  \[
    dt^2+\sum\delta_{ij}dx^idx^j
    =\sin^2(r/R)\,dr^2+\cos^2(r/R)\,dr^2\sum(s^i)^2+R^2\sin^2(r/R)\sum(ds^i)^2
  \]
  \[
    \qquad+2R\cos(r/R)\sin(r/R)\,dr\sum s^ids^i
    =dr^2+R^2\sin^2(r/R)\,ds_{n-1}^2,
  \]
  so $G$ pulls back the ambient metric on $\mathbb{R}^{n+1}$ to
  $dr^2+R^2\sin^2(r/R)ds_{n-1}^2=dr^2+\mathrm{sn}_{1/R^2}^2(r)ds_{n-1}^2$, exhibiting
  $S^n_{1/R^2}$ as isometric to $S^n(R)$. For $k=-1/R^2$: the analogous map
  $(s,r)\mapsto(R\sinh(r/R)s,R\cosh(r/R))$ into $\mathbb{R}^{n,1}$ satisfies, with
  $g_{\mathbb{R}^{n,1}}=\delta_{ij}dx^idx^j-dt^2$,
  \[
    -dt^2+\sum\delta_{ij}dx^idx^j
    =-\sinh^2(r/R)dr^2+\cosh^2(r/R)dr^2\sum(s^i)^2+R^2\sinh^2(r/R)\sum(ds^i)^2
  \]
  \[
    =dr^2+R^2\sinh^2(r/R)\,ds_{n-1}^2,
  \]
  by the same computation with $\cosh,\sinh$ in place of $\cos,\sin$; this
  exhibits $S^n_{-1/R^2}$ as isometric to $H^n(R)\subset\mathbb{R}^{n,1}$
  (\cref{ex:pet-ch1-hyperbolic-space}). The case $k=0$ is
  \cref{ex:pet-ch1-euclidean-coords} in polar form.
\end{proof}

\subsection{Doubly Warped Products}

\begin{definition}[doubly warped product metric]
  \label{def:pet-ch1-doubly-warped}
  \lean{PetersenLib.doublyWarpedProductMetric}
  \leanok
  \uses{def:pet-ch1-rot-sym-general}
  On $I\times S^p\times S^q$, a \emph{doubly warped product metric} is
  $dt^2+\rho^2(t)ds_p^2+\phi^2(t)ds_q^2$. Nondegeneracy forces $\rho,\phi$ not to
  vanish simultaneously; when $\rho(t)=0$, the smoothness analysis reduces to
  \cref{thm:pet-ch1-smoothness-criterion} applied to $\rho$ (with $\phi$ as a
  bystander factor).
  \source{petersen:page-0041}
\end{definition}

\begin{proposition}[Proposition 1.4.7 — smoothness at $t=0$]
  \label{prop:pet-ch1-doubly-warped-smooth-zero}
  \lean{PetersenLib.doublyWarpedSmoothAtZero}
  \leanok
  \uses{def:pet-ch1-doubly-warped,thm:pet-ch1-smoothness-criterion}
  If $\rho,\phi:(0,b)\to(0,\infty)$ are smooth (both warping functions of a
  doubly warped metric are positive on the interior), the doubly warped
  metric is smooth at $t=0$ iff $\rho(0)=0$, $\rho^{(\mathrm{even})}(0)=0$,
  $\dot\rho(0)=1$, $\phi(0)>0$, and $\phi^{(\mathrm{odd})}(0)=0$; the topology
  near $t=0$ is then $\mathbb{R}^{p+1}\times S^q$. (Positivity of $\phi$ near
  $0$ is needed for the ``only if'' direction: the metric coefficient
  $\phi^2$ determines $\phi$ only up to sign.)
  \source{petersen:page-0041}
\end{proposition}
\begin{proof}
  \leanok
  \uses{prop:pet-ch1-doubly-warped-smooth-zero,thm:pet-ch1-smoothness-criterion}
  In the Cartesian coordinates of \cref{thm:pet-ch1-smoothness-criterion} on
  the blown-up factor $\mathbb{R}^{p+1}$, the metric near $t=0$ has the
  rotationally symmetric part with coefficients
  $F_1=\rho^2(t)/t^2$, $F_2=1/t^2-\rho^2(t)/t^4$, and the bystander part
  $F_3\,g_{S^q}$ with $F_3=\phi^2(t)$, $t=|x|$. Smoothness of $F_1,F_2$ with
  $F_1(0)>0$ is, by \cref{thm:pet-ch1-smoothness-criterion}, equivalent to
  $\rho(0)=0$, $\dot\rho(0)=1$, $\rho^{(2l)}(0)=0$. It remains to show that
  $F_3$ extends smoothly with $F_3(0)>0$ iff $\phi(0)>0$ and
  $\phi^{(\mathrm{odd})}(0)=0$.

  ($\Leftarrow$) Split $\phi=\phi_e+\phi_o$ into even and odd parts. All
  derivatives of $\phi_o$ vanish at $0$: the even-order ones because $\phi_o$
  is odd, the odd-order ones because those of $\phi$ (hypothesis) and of the
  even function $\phi_e$ vanish. So $\phi_o$ is flat at $0$; by Whitney's
  even-function theorem $\phi_e(t)=g(t^2)$ with $g$ smooth, and a flat
  function of $t$ is a smooth function of $t^2$ as well. Since
  $t^2=\sum(x^i)^2$ is smooth in $x$, $\phi(|x|)=\phi_e(|x|)+\phi_o(|x|)$ is
  smooth, and $F_3=\phi(|x|)^2$ is smooth with $F_3(0)=\phi(0)^2>0$.

  ($\Rightarrow$) Restrict along a ray through a unit vector $e$:
  $c(t)=F_3(te)$ is smooth near $0$, even (both signs of $t$ give
  $\phi(|t|)^2$), positive near $0$ by continuity from $c(0)=F_3(0)>0$, and
  $c(t)=\phi(t)^2$ for small $t>0$. Correcting $c$ outside a small interval
  by an even cutoff makes it globally smooth, even and positive, so
  $\sqrt{c}$ is a globally smooth even function agreeing with $\phi$ near
  $0^+$ (positivity of $\phi$ fixes the sign of the root). Two smooth
  functions agreeing on an interval $(0,\delta)$ have equal derivatives at
  $0$, so all odd derivatives of $\phi$ vanish at $0$ by evenness of
  $\sqrt{c}$; moreover $\phi(0)^2=c(0)>0$ and $\phi(0)\ge0$ by
  right-continuity and positivity of $\phi$, so $\phi(0)>0$.
\end{proof}

\begin{proposition}[Proposition 1.4.8 — smoothness at $t=b$]
  \label{prop:pet-ch1-doubly-warped-smooth-b}
  \lean{PetersenLib.doublyWarpedSmoothAtB}
  \leanok
  \uses{def:pet-ch1-doubly-warped,thm:pet-ch1-smoothness-criterion}
  If $\rho,\phi:(0,b)\to(0,\infty)$ are smooth, the doubly warped
  metric is smooth at $t=b$ iff $\rho(b)=0$, $\rho^{(\mathrm{even})}(b)=0$,
  $\dot\rho(b)=-1$, $\phi(b)>0$, and $\phi^{(\mathrm{odd})}(b)=0$; the
  topology near $t=b$ is again $\mathbb{R}^{p+1}\times S^q$.
  \source{petersen:page-0041}
\end{proposition}
\begin{proof}
  \leanok
  \uses{prop:pet-ch1-doubly-warped-smooth-b,prop:pet-ch1-doubly-warped-smooth-zero}
  In the reflected variable $s=b-t$ the metric near $t=b$ takes the form
  treated in \cref{prop:pet-ch1-doubly-warped-smooth-zero}, with warping
  functions $\rho(b-s),\phi(b-s)$, positive for $s\in(0,b)$. Under the
  reflection $f\mapsto f(b-\cdot)$ a derivative of order $m$ picks up the
  sign $(-1)^m$, so even-order derivatives at the endpoint are preserved and
  odd-order ones change sign. The conditions of
  \cref{prop:pet-ch1-doubly-warped-smooth-zero} at $s=0$ --- value $0$,
  slope $+1$ and vanishing higher even-order derivatives for $\rho(b-\cdot)$;
  positive value and vanishing odd-order derivatives for $\phi(b-\cdot)$ ---
  therefore translate exactly to $\rho(b)=0$, $\dot\rho(b)=-1$,
  $\rho^{(2l)}(b)=0$, $\phi(b)>0$, $\phi^{(\mathrm{odd})}(b)=0$.
\end{proof}

\begin{remark}[topology types of doubly warped products]
  \label{rem:pet-ch1-doubly-warped-topologies}
  \lean{PetersenLib.doublyWarpedTopologyTypes}
  \leanok
  \uses{prop:pet-ch1-doubly-warped-smooth-zero,prop:pet-ch1-doubly-warped-smooth-b}
  If $\rho,\phi:[0,\infty)\to[0,\infty)$ are positive on $(0,\infty)$ and satisfy
  \cref{prop:pet-ch1-doubly-warped-smooth-zero}, the metric is smooth on
  $\mathbb{R}^{p+1}\times S^q$. If $\rho,\phi:[0,b]\to[0,\infty)$ are positive on $(0,b)$
  and satisfy both \cref{prop:pet-ch1-doubly-warped-smooth-zero,prop:pet-ch1-doubly-warped-smooth-b},
  the metric is smooth on $S^{p+1}\times S^q$. If instead the roles of $\rho,\phi$
  are interchanged at $t=b$, the metric is smooth on $S^{p+q+1}$.
  \source{petersen:page-0041}
\end{remark}

\begin{example}[Example 1.4.9 — spheres as doubly warped products]
  \label{ex:pet-ch1-sphere-doubly-warped}
  \lean{PetersenLib.sphereAsDoublyWarpedProduct}
  \leanok
  \uses{rem:pet-ch1-doubly-warped-topologies,ex:pet-ch1-sphere}
  $dt^2+\sin^2(t)ds_p^2+\cos^2(t)ds_q^2$, $t\in[0,\pi/2]$, is isometric to
  $(S^{p+q+1}(1),g_{S^{p+q+1}})$.
  \source{petersen:page-0041,petersen:page-0042}
\end{example}
\begin{proof}
  \leanok
  \uses{ex:pet-ch1-sphere-doubly-warped}
  Map $(0,\pi/2)\times S^p\times S^q\to\mathbb{R}^{p+1}\times\mathbb{R}^{q+1}$ by
  $(t,x,y)\mapsto(x\sin t,y\cos t)$, $|x|=|y|=1$: the image lies in the unit
  sphere of $\mathbb{R}^{p+q+2}$, and a computation identical in structure to the one
  in the proof of \cref{ex:pet-ch1-space-forms} (with $\rho=\sin$, $\phi=\cos$
  playing the roles of the two warping functions on the two factors) shows the
  pullback of the ambient Euclidean metric is
  $dt^2+\sin^2(t)ds_p^2+\cos^2(t)ds_q^2$.
\end{proof}

\subsection{Hopf Fibrations}

\begin{example}[Example 1.4.10 — the Hopf fibration revisited]
  \label{ex:pet-ch1-hopf-revisited}
  \lean{PetersenLib.hopfFibrationRevisitedSphere,PetersenLib.pullbackForm_sphereThreeParam,PetersenLib.pullbackForm_sphereTwoParam,PetersenLib.hopfMap_sphereThreeParam,PetersenLib.hopfFibrationRevisited}
  \leanok
  \uses{ex:pet-ch1-hopf-fibration,ex:pet-ch1-sphere-doubly-warped,ex:pet-ch1-su2-berger}
  On $S^3(1)$ write the metric as $dt^2+\sin^2(t)d\theta_1^2+\cos^2(t)d\theta_2^2$,
  $t\in[0,\pi/2]$, via
  $(t,e^{i\theta_1},e^{i\theta_2})\mapsto(\sin(t)e^{i\theta_1},\cos(t)e^{i\theta_2})
  \in S^3(1)\subset\mathbb{C}^2$; the Hopf fibres are
  $\theta\mapsto(t,e^{i(\theta_1+\theta)},e^{i(\theta_2+\theta)})$. On
  $S^2(1/2)$ use $dr^2+\tfrac14\sin^2(2r)d\theta^2$ via
  $(r,e^{i\theta})\mapsto(\tfrac12\cos(2r),\tfrac12\sin(2r)e^{i\theta})$. In
  these coordinates the Hopf map is
  $(t,e^{i\theta_1},e^{i\theta_2})\mapsto(t,e^{i(\theta_1-\theta_2)})$, agreeing
  with Wilhelm's map (\cref{ex:pet-ch1-hopf-fibration}); it is a Riemannian
  submersion, and composed with the identification
  $(t,e^{i\theta_1},e^{i\theta_2})\mapsto\begin{bsmallmatrix}\cos(t)e^{i\theta_1}&\sin(t)e^{i\theta_2}\\-\sin(t)e^{-i\theta_2}&\cos(t)e^{-i\theta_1}\end{bsmallmatrix}$
  gives an isometry $S^3(1)\to SU(2)$ with its left-invariant metric
  (\cref{ex:pet-ch1-su2-berger}).
  \source{petersen:page-0042,petersen:page-0043}
\end{example}
\begin{proof}
  \uses{ex:pet-ch1-hopf-revisited}
  \leanok
  On $S^3(1)$, $\partial_{\theta_1}+\partial_{\theta_2}$ is tangent to the Hopf
  fibre, $\partial_t$ and
  $u:=\big(\cos^2(t)\partial_{\theta_1}-\sin^2(t)\partial_{\theta_2}\big)/(\cos t\sin t)$
  are two further orthogonal unit vectors; on $S^2(1/2)$,
  $\partial_r,\ (2/\sin(2r))\partial_\theta$ is an orthonormal frame. Since the
  Hopf map sends $(t,\theta_1,\theta_2)\mapsto(r,\theta)=(t,\theta_1-\theta_2)$
  in these coordinates, $\partial_t\mapsto\partial_r$ and
  $u\mapsto\big(\cos^2(t)+\sin^2(t)\big)\partial_\theta/(\cos t\sin t)
  =(2/\sin(2r))\partial_\theta$, matching the two target basis vectors exactly;
  this shows the map is a linear isometry on the orthogonal complement of the
  fibre, i.e.\ a Riemannian submersion by \cref{def:pet-ch1-riemannian-submersion}.
\end{proof}

\begin{example}[Example 1.4.11 — the general submersion $I\times S^1\times S^1\to I\times S^1$]
  \label{ex:pet-ch1-hopf-general-submersion}
  \lean{PetersenLib.hopfFibrationGeneralSubmersion}
  \leanok
  \uses{def:pet-ch1-riemannian-submersion,def:pet-ch1-doubly-warped}
  The map $(t,e^{i\theta_1},e^{i\theta_2})\mapsto(t,e^{i(\theta_1-\theta_2)})$,
  $I\times S^1\times S^1\to I\times S^1$, is always a Riemannian submersion for
  the doubly warped domain metric $dt^2+\rho^2(t)d\theta_1^2+\phi^2(t)d\theta_2^2$
  and the rotationally symmetric target metric
  $dr^2+\dfrac{(\rho(t)\phi(t))^2}{\rho^2(t)+\phi^2(t)}\,d\theta^2$.
  \source{petersen:page-0043,petersen:page-0044}
\end{example}

\begin{example}[Example 1.4.12 — higher-dimensional Hopf submersion]
  \label{ex:pet-ch1-hopf-higher-dim}
  \lean{PetersenLib.hopfFibrationHigherDimSphere}
  \leanok
  \uses{ex:pet-ch1-hopf-general-submersion}
  On $I\times S^{2n+1}\times S^1$ with doubly warped metric
  $dt^2+\rho^2(t)ds_{2n+1}^2+\phi^2(t)d\theta^2$, split
  $ds_{2n+1}^2=h+g$ into the component $h$ along the Hopf fibre and the
  orthogonal component $g$ (via the $ds_{2n+1}^2$-orthogonal projection $\mathrm{pr}$
  onto the Hopf distribution: $h(v,w)=ds_{2n+1}^2(\mathrm{pr}(v),\mathrm{pr}(w))$,
  $g(v,w)=ds_{2n+1}^2(v-\mathrm{pr}(v),w-\mathrm{pr}(w))$). The unit circle acts
  isometrically and freely by simultaneous complex scalar multiplication on
  $S^{2n+1}\times S^1$; the quotient map onto
  $I\times((S^{2n+1}\times S^1)/S^1)=I\times S^{2n+1}$ is a Riemannian submersion
  for the target metric
  $dt^2+\rho^2(t)g+\dfrac{(\rho(t)\phi(t))^2}{\rho^2(t)+\phi^2(t)}h$.
  \source{petersen:page-0044}
\end{example}

\begin{example}[Example 1.4.13 — the case $n=1$]
  \label{ex:pet-ch1-hopf-su2-coframe}
  \lean{PetersenLib.hopfFibrationSUTwoCoframeSphere,PetersenLib.suTwoCoframe,PetersenLib.sphereMetricUnit_eq_sum_suTwoCoframe,PetersenLib.hopfComplement_eq_suTwoCoframe,PetersenLib.hopfQuotientForm_eq_suTwoCoframe}
  \leanok
  \uses{ex:pet-ch1-hopf-higher-dim,ex:pet-ch1-su2-berger}
  For $n=1$, using the coframe $\sigma^1,\sigma^2,\sigma^3$ from
  $S^3\simeq SU(2)$ (\cref{ex:pet-ch1-su2-berger}), the splitting
  $ds^2_{S^3}=h+g$ of \cref{ex:pet-ch1-hopf-higher-dim} reads
  $ds^2_{S^3}=(\sigma^1)^2+(\sigma^2)^2+(\sigma^3)^2$, $h=(\sigma^1)^2$,
  $g=(\sigma^2)^2+(\sigma^3)^2$, and the Riemannian submersion becomes
  \[
    \Big(I\times S^3\times S^1,\ dt^2+\rho^2(t)\big((\sigma^1)^2+(\sigma^2)^2+(\sigma^3)^2\big)+\phi^2(t)d\theta^2\Big)
  \]
  \[
    \longrightarrow
    \Big(I\times S^3,\ dt^2+\rho^2(t)\big((\sigma^2)^2+(\sigma^3)^2\big)+\tfrac{(\rho(t)\phi(t))^2}{\rho^2(t)+\phi^2(t)}(\sigma^1)^2\Big).
  \]
  \source{petersen:page-0045}
\end{example}

\begin{example}[Example 1.4.14 — the generalized Hopf fibration]
  \label{ex:pet-ch1-generalized-hopf}
  \lean{PetersenLib.generalizedHopfFibrationSphere,PetersenLib.generalizedHopfFibrationComplexProjective,PetersenLib.genHopfBase,PetersenLib.snocLpEquiv}
  \leanok
  \uses{ex:pet-ch1-hopf-higher-dim,ex:pet-ch1-fubini-study}
  Taking $\rho=\sin(t)$, $\phi=\cos(t)$, $t\in I=[0,\pi/2]$ in
  \cref{ex:pet-ch1-hopf-higher-dim} recovers the generalized Hopf fibration
  $S^{2n+3}\to\mathbb{CP}^{n+1}$ of \cref{ex:pet-ch1-fubini-study}, with the
  Fubini--Study metric represented as $dt^2+\sin^2(t)\big(g+\cos^2(t)h\big)$.
  \source{petersen:page-0045}
\end{example}

\section{Some Tensor Concepts}

\subsection{Type Change}

\begin{definition}[type change]
  \label{def:pet-ch1-tensor-type-change}
  \lean{PetersenLib.tensorTypeChange}
  \leanok
  \uses{def:pet-ch1-riemannian-manifold}
  For an $(s,t)$-tensor $T$, a section of
  $TM^{\otimes s}\otimes(T^*M)^{\otimes t}$, a Riemannian metric $g$ gives a
  natural isomorphism $TM\cong T^*M$, $v\mapsto(w\mapsto g(v,w))$, letting one
  replace occurrences of $TM$ by $T^*M$ or vice versa and so view $T$ as an
  $(s-k,t+k)$-tensor for any $k$ with $s-k,t+k\ge0$. Concretely, for a frame
  $E_1,\dots,E_n$ with coframe $\sigma^1,\dots,\sigma^n$, writing
  $T=T^{i_1\cdots i_s}_{j_1\cdots j_t}\,E_{i_1}\otimes\cdots\otimes E_{i_s}\otimes
  \sigma^{j_1}\otimes\cdots\otimes\sigma^{j_t}$, the dual of $E_i$ is
  $g(E_i,\cdot)=g_{ij}\sigma^j$ and the vector dual to $\sigma^j$ is $g^{ij}E_i$
  (where $(g^{ij})$ is the inverse matrix of $(g_{ij})$), so type change replaces
  each $E_i\leftrightarrow g_{ij}\sigma^j$ and $\sigma^j\leftrightarrow g^{ij}E_i$.
  For an orthonormal frame $g_{ij}=\delta_{ij}$ and indices are simply raised or
  lowered.
  \source{petersen:page-0045,petersen:page-0046,petersen:page-0047}
\end{definition}

\begin{notation}[the Ricci tensor as tensors of different types]
  \label{def:pet-ch1-ricci-types}
  \lean{PetersenLib.ricciTensorTypes}
  \leanok
  \uses{def:pet-ch1-tensor-type-change}
  As a $(1,1)$-tensor, $\mathrm{Ric}(E_i)=\mathrm{Ric}^j_iE_j$, i.e.\
  $\mathrm{Ric}=\mathrm{Ric}^j_i\,E_i\otimes\sigma^j$. As a $(0,2)$-tensor,
  $\mathrm{Ric}=\mathrm{Ric}_{jk}\,\sigma^j\otimes\sigma^k=g_{ji}\mathrm{Ric}^i_k\,
  \sigma^j\otimes\sigma^k$. As a $(2,0)$-tensor,
  $\mathrm{Ric}=\mathrm{Ric}^{ik}\,E_i\otimes E_k=g^{jl}\mathrm{Ric}^i_j\,E_i\otimes E_k$.
  \source{petersen:page-0047}
\end{notation}

\begin{notation}[the curvature tensor as tensors of different types]
  \label{def:pet-ch1-curvature-types}
  \lean{PetersenLib.curvatureTensorTypes,PetersenLib.curvatureTwoTwoTypes}
  \leanok
  \uses{def:pet-ch1-tensor-type-change}
  The $(1,3)$-curvature tensor $R(X,Y)Z$ is written
  $R=R^i_{jkl}\,E_i\otimes\sigma^j\otimes\sigma^l\otimes\sigma^k$. As a
  $(0,4)$-tensor, $R=R_{ijkl}\,\sigma^i\otimes\sigma^j\otimes\sigma^k\otimes\sigma^l
  =R^i_{jkl}g_{il}\cdot\sigma^i\otimes\sigma^j\otimes\sigma^k\otimes\sigma^l$ (with
  the raised index $l$ placed last, by convention). Raising the last index of the
  $(0,4)$-tensor produces the $(2,2)$-tensor
  $R=R^{kl}_{ij}\,E_k\otimes E_l\otimes\sigma^i\otimes\sigma^j
  =R^l_{ij}\,g^{sk}\cdots$, called the \emph{curvature operator}; raising a
  different index gives a different $(2,2)$-tensor.
  \source{petersen:page-0047,petersen:page-0048}
\end{notation}

\subsection{Contractions}

\begin{definition}[contraction]
  \label{def:pet-ch1-contraction}
  \lean{PetersenLib.tensorContraction}
  \leanok
  \uses{def:pet-ch1-tensor-type-change}
  \emph{Contractions} are traces of tensors. For a $(1,1)$-tensor
  $T=T^i_j\,E_i\otimes\sigma^j$, $C(T)=\mathrm{tr}\,T=T^i_i$; for $T=X\otimes\omega$
  (a vector field tensored with a $1$-form), $C(X\otimes\omega)=\omega(X)$, and
  contraction of a general tensor is a sum of such evaluations. For a $(0,2)$-tensor
  $T$, first type-change to a $(1,1)$-tensor and then trace:
  $C(T)=C(T_{ij}\sigma^i\otimes\sigma^j)=C(T_{ik}g^{kj}E_k\otimes\sigma^j)=T_{ik}g^{ki}$.
  \source{petersen:page-0048}
\end{definition}

\begin{remark}[the Ricci tensor as a contraction of curvature]
  \label{rem:pet-ch1-ricci-contraction}
  \lean{PetersenLib.ricciAsContraction}
  \leanok
  \uses{def:pet-ch1-contraction,def:pet-ch1-curvature-types,def:pet-ch1-ricci-types}
  $\mathrm{Ric}=\mathrm{Ric}^i_j\,E_i\otimes\sigma^j=R^{ij}_{ik}\,E_i\otimes\sigma^j
  =R^i_{iks}g^{sk}\,E_i\otimes\sigma^j$, equivalently as a $(0,2)$-tensor
  $\mathrm{Ric}=\mathrm{Ric}_{ij}\,\sigma^i\otimes\sigma^j=g^{kl}R_{kijl}\,\sigma^i\otimes\sigma^j$;
  after type change the two expressions agree.
  \source{petersen:page-0048}
\end{remark}

\begin{definition}[scalar curvature]
  \label{def:pet-ch1-scalar-curvature}
  \lean{PetersenLib.scalarCurvatureOfRicci}
  \leanok
  \uses{def:pet-ch1-contraction,def:pet-ch1-ricci-types}
  The \emph{scalar curvature} is the contraction of the Ricci tensor:
  \[
    \mathrm{scal}=\mathrm{tr}(\mathrm{Ric})=\mathrm{Ric}^i_i=R^i_{iks}g^{sk}
    =\mathrm{Ric}_{ik}g^{ki}=R_{ijkl}g^{ik}g^{jl}.
  \]
  \source{petersen:page-0048,petersen:page-0049}
\end{definition}

\subsection{Inner Products of Tensors}

\begin{definition}[operator norm and Euclidean norm]
  \label{def:pet-ch1-operator-norm}
  \lean{PetersenLib.operatorNorm}
  \leanok
  For a linear map $L:V\to W$ between normed spaces, the \emph{operator norm} is
  $\|L\|=\sup_{|v|=1}|Lv|$. If $L:V\to V$ is self-adjoint with eigenvalues
  $\lambda_1\le\dots\le\lambda_n$ (with multiplicity), $\|L\|=\max\{|\lambda_1|,|\lambda_n|\}$.
  \source{petersen:page-0049}
\end{definition}

\begin{definition}[Euclidean norm and trace inner product]
  \label{def:pet-ch1-euclidean-norm}
  \lean{PetersenLib.euclideanNormLinearMap,PetersenLib.traceInnerProduct}
  \leanok
  \uses{def:pet-ch1-contraction}
  The \emph{Euclidean norm} of $L:V\to W$ is
  $|L|=\sqrt{\mathrm{tr}(L^*\circ L)}=\sqrt{\mathrm{tr}(L\circ L^*)}$, where
  $L^*:W\to V$ is the adjoint; for self-adjoint $L:V\to V$ with eigenvalues
  $\lambda_i$, $|L|=\sqrt{\sum_i\lambda_i^2}$. It comes from the inner product
  $\langle L_1,L_2\rangle=\mathrm{tr}(L_1\circ L_2^*)=\mathrm{tr}(L_2\circ L_1^*)$.
  This Euclidean norm, not the operator norm, is used throughout.
  \source{petersen:page-0049}
\end{definition}

\begin{proposition}[Euclidean norm of a $(1,1)$-tensor]
  \label{prop:pet-ch1-norm-11-tensor}
  \lean{PetersenLib.euclideanNormOneOneTensor}
  \leanok
  \uses{def:pet-ch1-euclidean-norm,def:pet-ch1-tensor-type-change}
  For $T=T^i_j\,E_i\otimes\sigma^j$, viewed as a linear map $TM\to TM$,
  $|T|^2=T^j_iT^k_ig^{jl}g_{ki}$; for an orthonormal frame,
  $|T|^2=T^i_jT^j_i$.
  \source{petersen:page-0049,petersen:page-0050}
\end{proposition}
\begin{proof}
  \uses{prop:pet-ch1-norm-11-tensor}
  \leanok
  The dual map $T^*:T^*M\to T^*M$ is $T^*=T^j_i\,\sigma^i\otimes E_j$; type
  changing back to $TM\to TM$ via
  \cref{def:pet-ch1-tensor-type-change} gives
  $T^*=T^k_ig^{jl}g_{ki}\,E_j\otimes\sigma^i$. By
  \cref{def:pet-ch1-euclidean-norm}, $|T|^2=\mathrm{tr}(T\circ T^*)$, which in
  these components is $T^i_jT^k_ig^{jl}g_{ki}$ (contracting the free indices of
  $T$ against $T^*$ as in \cref{def:pet-ch1-contraction}); when the frame is
  orthonormal $g_{ij}=g^{ij}=\delta_{ij}$ and this collapses to $T^i_jT^j_i$.
\end{proof}

\begin{proposition}[Euclidean norm of a $(0,2)$-tensor]
  \label{prop:pet-ch1-norm-02-tensor}
  \lean{PetersenLib.euclideanNormZeroTwoTensor}
  \leanok
  \uses{prop:pet-ch1-norm-11-tensor}
  For $T=T_{ij}\,\sigma^i\otimes\sigma^j$, after type-changing to a $(1,1)$-tensor
  and applying \cref{prop:pet-ch1-norm-11-tensor}, $|T|^2=T_{ij}T^{ij}$.
  \source{petersen:page-0050}
\end{proposition}

\begin{definition}[pointwise inner product of tensors]
  \label{def:pet-ch1-pointwise-tensor-inner}
  \lean{PetersenLib.pointwiseTensorInnerProduct}
  \leanok
  \uses{def:pet-ch1-euclidean-norm}
  For $(s,t)$-tensors, declare that if $E_i$ is an orthonormal frame with dual
  coframe $\sigma^i$, the tensors $E_{i_1}\otimes\cdots\otimes E_{i_s}\otimes
  \sigma^{j_1}\otimes\cdots\otimes\sigma^{j_t}$ form an orthonormal basis; this
  determines the \emph{pointwise inner product} of two $(s,t)$-tensors,
  generalizing $g(X,Y)$ for vector fields.
  \source{petersen:page-0050}
\end{definition}

\begin{definition}[$L^2$ inner product of tensors]
  \label{def:pet-ch1-l2-tensor-inner}
  \lean{PetersenLib.l2TensorInnerProduct}
  \leanok
  \uses{def:pet-ch1-pointwise-tensor-inner,def:pet-ch1-volume-form-oriented}
  Integrating the pointwise inner product of two compactly supported tensors of
  the same type against the volume form gives an inner product on the space of
  such tensors: $(T_1,T_2)=\int_Mg(T_1,T_2)\,\mathrm{vol}$.
  \source{petersen:page-0050}
\end{definition}

\subsection{Positional Notation}

\begin{remark}[positional index notation]
  \label{rem:pet-ch1-positional-notation}
  \lean{PetersenLib.positionalTensorNotation}
  \leanok
  \uses{def:pet-ch1-tensor-type-change,prop:pet-ch1-norm-11-tensor}
  The frame notation for a section of $TM\otimes T^*M$ and one of
  $T^*M\otimes TM$ can look identical, causing confusion when computing
  Euclidean norms — $\sigma\otimes E$ and $E\otimes\sigma$ differ only in
  ordering, not in the map they represent, yet either could be read as a map
  $TM\to TM$. Writing indices positionally, e.g.\ a section
  $T$ of $T^*M\otimes TM\otimes T^*M$ as $T={T_i}^j{}_k\,\sigma^i\otimes
  E_j\otimes\sigma^k$, keeps the ordering unambiguous: the standard
  $(1,1)$-tensor is $T=T^i_j\,E_i\otimes\sigma^j$, while its adjoint (before type
  change), a section of $T^*M\otimes TM$, is
  $T^*=T^l_k\,\sigma^k\otimes E_l=T^l_jg_{kl}g^{ij}\,\sigma^k\otimes E_l$,
  recovering $|T|^2=T^i_jT^j_i$ as in \cref{prop:pet-ch1-norm-11-tensor}. This
  notation is not used consistently in the literature and is often replaced by
  invariant (coordinate-free) notation, but frame or coordinate computations
  remain necessary in both abstract and concrete settings.
  \source{petersen:page-0050,petersen:page-0051}
\end{remark}

\section{Exercises}

\begin{remark}[Exercise 1.6.1 — Cartesian product metrics]
  \label{rem:pet-ch1-ex-1}
  \lean{PetersenLib.exercise1_6_1}
  \leanok
  \uses{def:pet-ch1-riemannian-manifold,ex:pet-ch1-flat-torus}
  On $M\times N$ with product metric $g=g_M+g_N$: (1) show
  $(\mathbb{R}^n,g_{\mathbb{R}^n})=(\mathbb{R},dt^2)\times\cdots\times(\mathbb{R},dt^2)$ ($n$ factors);
  (2) show the flat square torus is
  $T^2=\mathbb{R}^2/\mathbb{Z}^2=\big(S^1,(1/2\pi)^2d\theta^2\big)\times\big(S^1,(1/2\pi)^2d\theta^2\big)$;
  (3) show $F(\theta_1,\theta_2)=\tfrac1{2\pi}(\cos\theta_1,\sin\theta_1,\cos\theta_2,\sin\theta_2)$
  is a Riemannian embedding $T^2\to\mathbb{R}^4$.
  \source{petersen:page-0051,petersen:page-0052}
\end{remark}

\begin{remark}[Exercise 1.6.2 — quotient metric for an isometric action]
  \label{rem:pet-ch1-ex-2}
  \lean{PetersenLib.exercise1_6_2}
  \leanok
  \uses{def:pet-ch1-riemannian-submersion}
  If a group $G$ acts isometrically on $(M,g)$ so that $M/G$ is a manifold and
  the quotient map a submersion, show there is a unique Riemannian metric on
  $M/G$ making the quotient map a Riemannian submersion.
  \source{petersen:page-0052}
\end{remark}

\begin{remark}[Exercise 1.6.3 — volume of a finite cover]
  \label{rem:pet-ch1-ex-3}
  \lean{PetersenLib.exercise1_6_3}
  \leanok
  \uses{rem:pet-ch1-covering-metric,def:pet-ch1-volume-form-oriented}
  For a Riemannian $k$-fold covering map $M\to N$ ($k\ge1$, $N$ second
  countable, Riemannian measures Borel), show
  $\mathrm{vol}\,M=k\cdot\mathrm{vol}\,N$.
  \source{petersen:page-0052}
\end{remark}
\begin{proof}
  \leanok
  \uses{rem:pet-ch1-ex-3}
  Each point of $N$ has an evenly covered open neighbourhood; by second
  countability (Lindel\"of) countably many such sets $U_0,U_1,\dots$ cover
  $N$, and disjointification $B_i=U_i\setminus(U_0\cup\dots\cup U_{i-1})$
  turns them into countably many disjoint Borel pieces, each still inside an
  evenly covered set. Over a piece $B\subseteq U$ the trivialization
  $F^{-1}(U)\cong U\times F^{-1}(y)$ splits $F^{-1}(B)$ into $k$ disjoint
  sheets, each open in $F^{-1}(U)$ intersected with the Borel set
  $F^{-1}(B)$, on each of which $F$ is injective with image exactly $B$;
  since the covering is a local isometry, $F$ is measure preserving on each
  sheet, so $\mathrm{vol}\,F^{-1}(B)=k\cdot\mathrm{vol}\,B$. Summing over
  the pieces by countable additivity gives
  $\mathrm{vol}\,M=k\cdot\mathrm{vol}\,N$. (For $k=0$ the hypothesis
  ``$\#F^{-1}(y)=0$'' would also admit infinite fibers, e.g.
  $\mathbb{R}\to S^1$, for which the statement fails; hence $k\ge1$.)
\end{proof}

\begin{remark}[Exercise 1.6.4 — volume form of a warped product]
  \label{rem:pet-ch1-ex-4}
  \lean{PetersenLib.exercise1_6_4}
  \leanok
  \uses{def:pet-ch1-doubly-warped,def:pet-ch1-volume-form-local}
  Show the volume form of $dr^2+\rho^2(r)g_N$ on $I\times N$ is
  $\rho^{n-1}\,dr\wedge\mathrm{vol}_N$.
  \source{petersen:page-0052}
\end{remark}

\begin{remark}[Exercise 1.6.5 — dual frame and the volume form]
  \label{rem:pet-ch1-ex-5}
  \lean{PetersenLib.exercise1_6_5}
  \leanok
  \uses{def:pet-ch1-volume-form-local}
  For an orthonormal frame $E_1,\dots,E_n$, show the dual coframe is
  $\sigma^i(X)=g(E_i,X)$ and $\mathrm{vol}=\pm\sigma^1\wedge\cdots\wedge\sigma^n$.
  \source{petersen:page-0052}
\end{remark}

\begin{remark}[Exercise 1.6.6 — volume form in local coordinates]
  \label{rem:pet-ch1-ex-6}
  \lean{PetersenLib.exercise1_6_6}
  \leanok
  \uses{def:pet-ch1-volume-form-local,def:pet-ch1-coordinate-metric}
  Show that in local coordinates $x^1,\dots,x^n$,
  $\mathrm{vol}=\pm\sqrt{\det[g_{ij}]}\,dx^1\wedge\cdots\wedge dx^n$.
  \source{petersen:page-0052}
\end{remark}

\begin{remark}[Exercise 1.6.7 — paper models of warped products]
  \label{rem:pet-ch1-ex-7}
  \lean{PetersenLib.exercise1_6_7}
  \leanok
  \uses{rem:pet-ch1-rotationally-symmetric}
  Construct paper models of $dt^2+a^2t^2d\theta^2$: for $a=1$ this is the
  Euclidean plane, for $a<1$ a cone; describe the case $a>1$.
  \source{petersen:page-0052}
\end{remark}

\begin{remark}[Exercise 1.6.8 — smoothness for radius $R$]
  \label{rem:pet-ch1-ex-8}
  \lean{PetersenLib.exercise1_6_8}
  \leanok
  \uses{thm:pet-ch1-smoothness-criterion}
  For $dr^2+\rho^2(r)g_{S^{n-1}(R)}$ with $\rho(0)=0$, show smoothness at $r=0$
  requires $\dot\rho(0)=1/R$ and $\rho^{(2k)}(0)=0$.
  \source{petersen:page-0052}
\end{remark}

\begin{remark}[Exercise 1.6.9 — $\mathrm{Iso}(\mathbb{R}^n)$ as a matrix group]
  \label{rem:pet-ch1-ex-9}
  \lean{PetersenLib.exercise1_6_9}
  \leanok
  \uses{def:pet-ch1-isometry-group,ex:pet-ch1-isometry-group-euclidean}
  Viewing $\mathbb{R}^n$ as the hyperplane $x^{n+1}=R$ in $\mathbb{R}^{n+1}$, show
  $\mathrm{Iso}(\mathbb{R}^n)$ is identified with matrices
  $\begin{bsmallmatrix}O&v\\0&1\end{bsmallmatrix}$, $v\in\mathbb{R}^n$, $O\in O(n)$, and
  that these are exactly the linear maps preserving $x^{n+1}=R$ and the
  degenerate form $x^1y^1+\dots+x^ny^n$.
  \source{petersen:page-0052}
\end{remark}

\begin{remark}[Exercise 1.6.10 — Sylvester basis for a nondegenerate form]
  \label{rem:pet-ch1-ex-10}
  \lean{PetersenLib.exercise1_6_10}
  \leanok
  \uses{def:pet-ch1-pseudo-riemannian}
  For $V$ with a symmetric nondegenerate bilinear form $g$ of index $p$: (1) show
  there is a basis $e_1,\dots,e_n$ with $g(e_i,e_j)=0$ ($i\ne j$),
  $g(e_i,e_i)=1$ for $i\le n-p$, $g(e_i,e_i)=-1$ for $i>n-p$, so $V\cong\mathbb{R}^{p,q}$;
  (2) show $v=\sum_{i\le n-p}g(v,e_i)e_i-\sum_{i>n-p}g(v,e_i)e_i$; (3) for
  $L:V\to V$ linear, show
  $\mathrm{tr}(L)=\sum_ig(L(e_i),e_i)/g(e_i,e_i)$.
  \source{petersen:page-0052,petersen:page-0053}
\end{remark}

\begin{remark}[Exercise 1.6.11 — type change as contraction with $g^{-1}$]
  \label{rem:pet-ch1-ex-11}
  \lean{PetersenLib.exercise1_6_11}
  \leanok
  \uses{def:pet-ch1-tensor-type-change,def:pet-ch1-contraction}
  With $g^{-1}$ the $(2,0)$-tensor inner product on $T^*M$, show type change of a
  tensor is a contraction of its tensor product with $g$ or $g^{-1}$.
  \source{petersen:page-0053}
\end{remark}

\begin{remark}[Exercise 1.6.12 — norm via an orthonormal basis]
  \label{rem:pet-ch1-ex-12}
  \lean{PetersenLib.exercise1_6_12}
  \leanok
  \uses{def:pet-ch1-euclidean-norm}
  For a $(1,1)$-tensor $T$ and orthonormal basis $E_i$, show
  $|T|^2=\sum_i|T(E_i)|^2$.
  \source{petersen:page-0053}
\end{remark}

\begin{remark}[Exercise 1.6.13 — orthogonality of symmetric and skew tensors]
  \label{rem:pet-ch1-ex-13}
  \lean{PetersenLib.exercise1_6_13}
  \leanok
  \uses{def:pet-ch1-pointwise-tensor-inner}
  For $(1,1)$-tensors $S,T$ with $S$ symmetric and $T$ skew-symmetric, show
  $g(S,T)=0$.
  \source{petersen:page-0053}
\end{remark}

\begin{remark}[Exercise 1.6.14 — inner product via type change and contraction]
  \label{rem:pet-ch1-ex-14}
  \lean{PetersenLib.exercise1_6_14}
  \leanok
  \uses{def:pet-ch1-pointwise-tensor-inner,def:pet-ch1-tensor-type-change,def:pet-ch1-contraction}
  Show the inner product of two tensors of the same type can be described as
  (possibly several) type changes of one tensor followed by (possibly several)
  contractions.
  \source{petersen:page-0053}
\end{remark}

\begin{remark}[Exercise 1.6.15 — the quotient map $\mathbb{F}^{n+1}-\{0\}\to\mathbb{FP}^n$]
  \label{rem:pet-ch1-ex-15}
  \lean{PetersenLib.exercise1_6_15}
  \leanok
  \uses{def:pet-ch1-riemannian-submersion,ex:pet-ch1-fubini-study}
  For $F(x)=\mathrm{span}_{\mathbb{F}}\{x\}$, $\mathbb{F}=\mathbb{R},\mathbb{C}$, with $\mathbb{FP}^n$ metrized so
  that $F$ restricted to the unit sphere is a Riemannian submersion: (1) show $F$
  is a submersion ($F$ factors as the Hopf submersion after the radial
  projection $x\mapsto x/|x|$, both of which have surjective differentials);
  (2) show $F$ is not a Riemannian submersion for the standard
  metric on $\mathbb{F}^{n+1}-\{0\}$ (by scale invariance, horizontal vectors at $x$
  and $2x$ would both map isometrically, forcing them to vanish); (3) determine
  whether some metric on $\mathbb{F}^{n+1}-\{0\}$ makes $F$ a Riemannian submersion —
  yes: the cylinder metric $dr^2+g_{S^n}$, i.e.\ $g_x(u,v)=\langle u_r,v_r\rangle
  +|x|^{-2}\langle u_T,v_T\rangle$ in the radial/tangential decomposition, whose
  horizontal vectors at $x$ are the Euclidean-horizontal ones scaled by $|x|$.
  \source{petersen:page-0053}
\end{remark}

\begin{remark}[Exercise 1.6.16 — arc length]
  \label{rem:pet-ch1-ex-16}
  \lean{PetersenLib.exercise1_6_16}
  \leanok
  \uses{def:pet-ch1-riemannian-manifold}
  For $c:[a,b]\to(M,g)$, $L(c)=\int_{[a,b]}|\dot c|\,dt$: (1) show $L(c)$ is
  independent of parametrization; (2) show any curve of nowhere-vanishing speed
  reparametrizes to unit speed; (3) show arc length extends to absolutely
  continuous curves, and that absolute continuity is well defined for curves in
  manifolds.
  \source{petersen:page-0053,petersen:page-0054}
\end{remark}

\begin{remark}[Exercise 1.6.17 — arc length under Riemannian immersions]
  \label{rem:pet-ch1-ex-17}
  \lean{PetersenLib.exercise1_6_17}
  \leanok
  \uses{rem:pet-ch1-ex-16,def:pet-ch1-pullback-metric}
  Show arc length is preserved by Riemannian immersions.
  \source{petersen:page-0054}
\end{remark}

\begin{remark}[Exercise 1.6.18 — length under a Riemannian submersion]
  \label{rem:pet-ch1-ex-18}
  \lean{PetersenLib.exercise1_6_18}
  \leanok
  \uses{rem:pet-ch1-ex-16,def:pet-ch1-riemannian-submersion}
  For a Riemannian submersion $F:(M,g_M)\to(N,g_N)$ and $c:[a,b]\to M$, show
  $L(F\circ c)\le L(c)$, with equality iff $\dot c(t)\perp\ker DF_{c(t)}$ for all
  $t$.
  \source{petersen:page-0054}
\end{remark}

\begin{remark}[Exercise 1.6.19 — straight lines minimize in $\mathbb{R}^n$]
  \label{rem:pet-ch1-ex-19}
  \lean{PetersenLib.exercise1_6_19}
  \leanok
  \uses{rem:pet-ch1-ex-16,ex:pet-ch1-euclidean-space}
  Show directly that any curve between two points of Euclidean space is at least
  as long as the Euclidean distance between them, with equality only for a
  curve lying on the connecting segment. Hint: for an appropriate unit vector
  $v$, compare the length of $v\cdot c(t)$ to the length of $c$.
  \source{petersen:page-0054}
\end{remark}

\begin{proposition}[No flat isometric immersion into an umbilic hypersurface]
  \label{prop:pet-ch1-no-flat-immersion}
  \lean{PetersenLib.no_isometricImmersion_flat_to_umbilic}
  \leanok
  \uses{def:pet-ch1-riemannian-manifold,def:pet-ch1-pullback-metric}
  Let $\beta$ be a symmetric bilinear form on an $(n+1)$-dimensional real
  vector space $E$, let $c\ne0$, and let $n\ge2$. Then there is no
  $C^\infty$ map $G:U\to E$, $U\subset\mathbb{R}^n$ open and nonempty, with
  \[
    \beta(G,G)=c,\qquad \beta(\partial_iG,\partial_jG)=\delta_{ij}
    \quad\text{on }U .
  \]
  In words: an open piece of flat $\mathbb{R}^n$ admits no isometric
  immersion into the level set $\{\beta(\cdot,\cdot)=c\}$ of $\beta$, of which
  $S^n\subset\mathbb{R}^{n+1}$ ($\beta=\langle\cdot,\cdot\rangle$, $c=1$) and
  $H^n\subset\mathbb{R}^{n,1}$ ($\beta=\eta$, $c=-1$) are the two instances.
\end{proposition}

\begin{proof}
  \leanok
  Differentiating $\beta(G,G)=c$ gives $\beta(\partial_iG,G)=0$: every first
  derivative is $\beta$-orthogonal to the normal $G$. Differentiating that
  relation again gives
  \[
    \beta(\partial_i\partial_jG,G)=-\beta(\partial_jG,\partial_iG)=-\delta_{ij}.
  \]
  Next put $A_{ij,k}:=\beta(\partial_i\partial_jG,\partial_kG)$. Differentiating
  $\beta(\partial_jG,\partial_kG)=\delta_{jk}$ in the direction $i$ gives
  $A_{ij,k}=-A_{ik,j}$, while equality of mixed partials gives
  $A_{ij,k}=A_{ji,k}$. Chaining the two,
  \[
    A_{ij,k}=-A_{ik,j}=-A_{ki,j}=A_{kj,i}=A_{jk,i}=-A_{ji,k}=-A_{ij,k},
  \]
  so $A_{ij,k}=0$: the second derivatives are $\beta$-orthogonal to all
  tangent directions.

  The $n+1$ vectors $\partial_1G,\dots,\partial_nG,G$ are $\beta$-orthogonal
  with $\beta$-squares $1,\dots,1,c$, all nonzero, hence linearly independent,
  hence a basis of $E$ (this is where $\dim E=n+1$ is used). Expanding
  $\partial_i\partial_jG$ in that basis, the two displays above compute all its
  coefficients:
  \[
    \partial_i\partial_jG=-\frac{\delta_{ij}}{c}\,G .
  \]
  Finally fix two distinct indices $1,2$ (here $n\ge2$ enters) and consider
  $\Phi:=\partial_2G$. By the last display $\partial_v\Phi=-c^{-1}\langle
  e_2,v\rangle\,G$, so differentiating once more,
  \[
    \partial_w\partial_v\Phi=-c^{-1}\langle e_2,v\rangle\,\partial_wG .
  \]
  Symmetry of this second derivative in $(v,w)$, evaluated at $v=e_2$, $w=e_1$,
  yields $-c^{-1}\partial_1G=0$, whence $\partial_1G=0$ since $c\ne0$. But
  $\beta(\partial_1G,\partial_1G)=1\ne0$, a contradiction.
\end{proof}

\begin{remark}[Exercise 1.6.20 — great circles minimize on the sphere]
  \label{rem:pet-ch1-ex-20}
  \lean{PetersenLib.exercise1_6_20}
  \leanok
  \uses{ex:pet-ch1-sphere,rem:pet-ch1-ex-16,prop:pet-ch1-no-flat-immersion}
  For $p,q\in S^n\subset\mathbb{R}^{n+1}$, $v\in T_pS^n$ unit: (1) show $p\cos t+v\sin t$
  is a unit-speed curve through $p$ with velocity $v$ at $t=0$; (2) show
  $F(r,v)=p\cos r+v\sin r$, $r\in[0,\pi]$, $v\perp p$, $|v|=1$, is a
  diffeomorphism $(0,\pi)\times S^{n-1}\to S^n-\{\pm p\}$; (3) for the radial
  field $\partial_r=F_*(\partial_r)$ and $q=F(r_0,v_0)$, show
  $\partial_r|_q=\big(-p+(p\cdot q)q\big)/\sqrt{1-(p\cdot q)^2}=-p\sin r_0+v_0\cos r_0$;
  (4) show any curve from $p$ to $q$ is longer than $r_0$ unless it lies on the
  great circle, using $\dot c\cdot\partial_r=dr/dt$; (5) conclude there is no
  Riemannian immersion of an open $U\subset\mathbb{R}^n$ into $S^n$, using that
  equilateral spherical triangles bounded by great-circle arcs always have
  angles $>\pi/3$. (Part (5) is formalized through
  Proposition~\ref{prop:pet-ch1-no-flat-immersion} rather than through the
  angle comparison suggested here: the latter needs the image of a straight
  segment under a local isometry to be a great-circle arc, i.e.\ a geodesic of
  the target, which is a fact of Chapter~5.)
  \source{petersen:page-0054}
\end{remark}

\begin{remark}[Exercise 1.6.21 — hyperbolas minimize on hyperbolic space]
  \label{rem:pet-ch1-ex-21}
  \lean{PetersenLib.exercise1_6_21}
  \leanok
  \uses{ex:pet-ch1-hyperbolic-space,rem:pet-ch1-ex-16,prop:pet-ch1-no-flat-immersion}
  For $p,q\in H^n\subset\mathbb{R}^{n,1}$, $v\in T_pH^n$ unit ($|p|^2=|q|^2=-1$,
  $|v|^2=1$, $p\cdot v=0$): (1) show $p\cosh t+v\sinh t$ is a unit-speed curve
  through $p$ with velocity $v$; (2) show $F(r,v)=p\cosh r+v\sinh r$, $r\ge0$,
  $v\cdot p=0$, $|v|^2=1$, is a diffeomorphism $(0,\infty)\times S^{n-1}\to H^n-\{p\}$;
  (3) for $\partial_r=F_*(\partial_r)$ and $q=F(r_0,v_0)$, show
  $\partial_r|_q=\big(-p-(q\cdot p)q\big)/\sqrt{-1+(q\cdot p)^2}=p\sinh r_0+v_0\cosh r_0$;
  (4) show any curve from $p$ to $q$ is longer than $r_0$ unless it lies on the
  hyperbola; (5) conclude there is no Riemannian immersion of an open
  $U\subset\mathbb{R}^n$ into $H^n$, using that equilateral hyperbolic triangles
  bounded by hyperbola arcs always have angles $<\pi/3$. (As for
  Exercise~1.6.20, part (5) is formalized through
  Proposition~\ref{prop:pet-ch1-no-flat-immersion}, with $\beta=\eta$ and
  $c=-1$: the sign of the ambient form is irrelevant to the obstruction, so
  $H^n$ is as rigid as $S^n$ against flat immersions.)
  \source{petersen:page-0054,petersen:page-0055}
\end{remark}

\begin{remark}[Exercise 1.6.22 (F. Wilhelm) — the quaternionic Hopf fibrations]
  \label{rem:pet-ch1-ex-22}
  \lean{PetersenLib.exercise1_6_22}
  \leanok
  \uses{ex:pet-ch1-hopf-fibration}
  Writing quaternions $q=a+bi+cj+dk=z+wj$: (1) show $\mathbb{H}$ is realized as the
  matrix algebra $q=\begin{bsmallmatrix}z&w\\-\bar w&\bar z\end{bsmallmatrix}$ with
  $i,j,k$ the indicated matrices, giving $\mathbb{R}$-bilinear associative
  multiplication; (2) for $\bar q=a-bi-cj-dk$, show
  $a^2+b^2+c^2+d^2=|q|^2=q\bar q=\bar qq=|z|^2+|w|^2=\det\begin{bsmallmatrix}z&w\\-\bar w&\bar z\end{bsmallmatrix}$,
  $|pq|=|p||q|$, and $\overline{pq}=\bar q\bar p$; (3) for
  $H^l(p,q)=\big(\tfrac12(|p|^2-|q|^2),\bar pq\big)$ and
  $H^r(p,q)=\big(\tfrac12(|p|^2-|q|^2),p\bar q\big)$, show both map
  $S^7(1)\subset\mathbb{H}^2$ to $S^4(1/2)\subset\mathbb{R}\oplus\mathbb{H}$; (4) show the fibres of
  $H^l$ are the orbits of left multiplication by unit quaternions on $\mathbb{H}^2$;
  (5) show the fibres of $H^r$ are the orbits of right multiplication by unit
  quaternions; (6) show both $H^l,H^r:S^7(1)\to S^4(1/2)$ are Riemannian
  submersions.
  \source{petersen:page-0055,petersen:page-0056}
\end{remark}

\begin{remark}[Exercise 1.6.23 — Euler-number vector bundles over $S^2$]
  \label{rem:pet-ch1-ex-23}
  \lean{PetersenLib.exercise1_6_23}
  \leanok
  \uses{ex:pet-ch1-hopf-su2-coframe,thm:pet-ch1-smoothness-criterion}
  For the Riemannian submersion of \cref{ex:pet-ch1-hopf-su2-coframe} with
  $\rho,\phi>0$ on $(0,\infty)$, set $f=\rho$ and
  $h=(\rho(t)\phi(t))^2/(\rho^2(t)+\phi^2(t))$, and assume $f(0)>0$,
  $f^{(\mathrm{odd})}(0)=0$, $h(0)=0$, $h'(0)=k$, $h^{(\mathrm{even})}(0)=0$ for a
  positive integer $k$. Show this yields a smooth metric on the vector bundle
  over $S^2$ with Euler number $\pm k$ (away from the zero section this bundle
  is $(0,\infty)\times S^3/\mathbb{Z}_k$, realized as a submersion of $S^3\times\mathbb{R}^2$;
  $k=2$ gives $TS^2$, $k=1$ gives $\mathbb{RP}^2-\{\mathrm{pt}\}$).
  \source{petersen:page-0056,petersen:page-0057}
\end{remark}

\begin{remark}[Exercise 1.6.24 — bi-invariant metrics on compact Lie groups]
  \label{rem:pet-ch1-ex-24}
  \lean{PetersenLib.exercise1_6_24}
  \uses{def:pet-ch1-left-invariant-metric,def:pet-ch1-volume-form-local}
  Let $G$ be a compact Lie group. (1) Show $G$ admits a bi-invariant metric, by
  averaging a left-invariant metric $g_l$ over right translations against the
  left-invariant volume form:
  $g(v,w)=\big(\int_G\mathrm{vol}\big)^{-1}\int_Gg_l(DR_x(v),DR_x(w))\,\mathrm{vol}$;
  (2) show conjugation $\mathrm{Ad}_h(x)=hxh^{-1}$ is a Riemannian isometry for
  any bi-invariant metric, so its differential $\mathrm{Ad}_h:\mathfrak g\to\mathfrak g$
  at $e$ is a linear isometry; (3) deduce that
  $\mathrm{ad}_U:\mathfrak g\to\mathfrak g$, $\mathrm{ad}_UX=[U,X]$, is skew-symmetric:
  $g([U,X],Y)=-g(X,[U,Y])$.
  \source{petersen:page-0057}
\end{remark}

\begin{remark}[Exercise 1.6.25 — characterization of bi-invariant pseudo-metrics]
  \label{rem:pet-ch1-ex-25}
  \lean{PetersenLib.exercise1_6_25}
  \leanok
  \uses{def:pet-ch1-pseudo-riemannian,def:pet-ch1-left-invariant-metric}
  For a Lie group $G$ with Lie algebra $\mathfrak g$, show a nondegenerate,
  symmetric, bilinear form $(X,Y)$ on $\mathfrak g$ defines a bi-invariant
  pseudo-Riemannian metric iff $(X,Y)=(\mathrm{Ad}_hX,\mathrm{Ad}_hY)$ for all
  $h\in G$.
  \source{petersen:page-0057}
\end{remark}

\begin{remark}[Exercise 1.6.26 — averaging to a $G$-invariant metric]
  \label{rem:pet-ch1-ex-26}
  \lean{PetersenLib.exercise1_6_26}
  \uses{def:pet-ch1-riemannian-manifold}
  For a compact group $G$ acting on a manifold $M$, show $M$ admits a Riemannian
  metric for which $G$ acts by isometries. Hint: first show any manifold admits
  a Riemannian metric (via a partition of unity), then average it over $G$.
  \source{petersen:page-0058}
\end{remark}

\begin{remark}[Exercise 1.6.27 — the Killing form]
  \label{rem:pet-ch1-ex-27}
  \lean{PetersenLib.exercise1_6_27}
  \leanok
  \uses{def:pet-ch1-left-invariant-metric}
  For a Lie group $G$ with Killing form $B(X,Y)=\mathrm{tr}(\mathrm{ad}_X\circ\mathrm{ad}_Y)$
  on $\mathfrak g$: (1) show $B$ is symmetric and bilinear; (2) if $G$ admits a
  bi-invariant metric, show $B(X,X)\le0$; (3) show
  $B(\mathrm{ad}_ZX,Y)=-B(X,\mathrm{ad}_ZY)$; (4) for $G$ connected, show
  $B(\mathrm{Ad}_hX,\mathrm{Ad}_hY)=B(X,Y)$ by showing
  $t\mapsto B(\mathrm{Ad}_{\exp(tZ)}X,\mathrm{Ad}_{\exp(tZ)}Y)$ is constant. When
  $\mathfrak g$ is semisimple, $B$ is nondegenerate and (with a sign, $-B$, when
  $G$ is also compact) gives a bi-invariant pseudo-/Riemannian metric.
  \source{petersen:page-0058}
\end{remark}

\begin{remark}[Exercise 1.6.28 — $\mathrm{SL}(n,\mathbb{R})$]
  \label{rem:pet-ch1-ex-28}
  \lean{PetersenLib.exercise1_6_28}
  \leanok
  \uses{def:pet-ch1-pseudo-riemannian,def:pet-ch1-left-invariant-metric}
  For $\mathrm{SL}(n,\mathbb{R})$ with Lie algebra $\mathfrak{sl}(n,\mathbb{R})$ (trace-zero
  matrices), show the symmetric bilinear form $(X,Y)=\mathrm{tr}(XY)$ is
  nondegenerate and $\mathrm{Ad}_h$-invariant, hence defines a bi-invariant
  pseudo-Riemannian metric on $\mathrm{SL}(n,\mathbb{R})$.
  \source{petersen:page-0058}
\end{remark}

\begin{remark}[Exercise 1.6.29 — a group with no bi-invariant pseudo-metric]
  \label{rem:pet-ch1-ex-29}
  \lean{PetersenLib.exercise1_6_29}
  \leanok
  \uses{def:pet-ch1-pseudo-riemannian,def:pet-ch1-left-invariant-metric}
  Show the matrices $\mathrm{diag}(a^{-1},ab,1)$-type
  $\begin{bsmallmatrix}a^{-1}&0&0\\0&ab&\ast\\0&0&1\end{bsmallmatrix}$, $a>0$,
  $b\in\mathbb{R}$, define a $2$-dimensional Lie group admitting no bi-invariant
  pseudo-Riemannian metric.
  \source{petersen:page-0058}
\end{remark}
